% Public reading edition. Mathematical body retained from the author's manuscript.
\documentclass[11pt,a4paper]{article}
\usepackage[margin=1in,headheight=14pt,footskip=30pt]{geometry}
\usepackage{fontspec}
\setmainfont{lmroman10-regular.otf}[BoldFont=lmroman10-bold.otf,ItalicFont=lmroman10-italic.otf,SmallCapsFont=lmromancaps10-regular.otf]
\usepackage{amsmath,amsfonts,amssymb,mathtools,mathrsfs}
\usepackage[amsmath,thmmarks,hyperref]{ntheorem}
\usepackage{longtable,booktabs,array,calc,etoolbox,graphicx}
\usepackage{xcolor}
\usepackage[unicode=true,hypertexnames=false,colorlinks=true,linkcolor=blue,citecolor=blue,urlcolor=blue,
  bookmarksopen=false,bookmarksnumbered=false]{hyperref}
\urlstyle{same}
\newcommand{\proofbox}{\ensuremath{\square}}
\hypersetup{pdfauthor={Letao Ouyang},
  pdftitle={Critical Dependence Stress for Diversification Gains under Fixed Marginals},
  pdfsubject={Author manuscript - public reading edition},
  pdfcreator={XeLaTeX}}
\makeatletter
\renewtheoremstyle{plain}
 {\item[\hskip\labelsep\theorem@headerfont ##1\ \textup{##2\theorem@separator}]}
 {\item[\hskip\labelsep\theorem@headerfont ##1\ \textup{##2}]{\normalfont\bfseries\textup{(##3)\theorem@separator}}\ }
\renewtheoremstyle{nonumberplain}
 {\item[\hskip\labelsep\theorem@headerfont ##1\theorem@separator]}
 {\item[\hskip\labelsep\theorem@headerfont ##3\theorem@separator]}
\makeatother
\theoremstyle{nonumberplain}
\theoremheaderfont{\normalfont\itshape}
\theorembodyfont{\normalfont}
\theoremseparator{.}
\theoremsymbol{\ensuremath{\square}}
\newtheorem{proof}{Proof}
\theoremstyle{plain}
\theoremheaderfont{\normalfont\bfseries}
\theorembodyfont{\normalfont\itshape}
\theoremseparator{.}
\theoremsymbol{}
\newtheorem{theorem}{Theorem}
\newtheorem{lemma}[theorem]{Lemma}
\newtheorem{proposition}[theorem]{Proposition}
\newtheorem{corollary}[theorem]{Corollary}
\theorembodyfont{\normalfont}
\newtheorem{definition}[theorem]{Definition}
\newtheorem{remark}[theorem]{Remark}
\newcommand{\resultnumber}[1]{\renewcommand{\thetheorem}{#1}}
\newcommand{\keepproofstart}[1]{\par\begingroup
  \dimen0=\pagegoal\advance\dimen0 by-\pagetotal
  \ifdim\dimen0<#1\relax\newpage\fi\endgroup}
\newenvironment{resultparts}{%
  \begin{list}{(\roman{enumi})}{\usecounter{enumi}%
    \setlength{\leftmargin}{2em}\setlength{\labelwidth}{1.7em}%
    \setlength{\labelsep}{0.3em}\setlength{\itemsep}{0.4ex}}%
}{\end{list}}
\providecommand{\tightlist}{\setlength{\itemsep}{0pt}\setlength{\parskip}{0pt}}
\providecommand{\passthrough}[1]{#1}
\AtBeginEnvironment{thebibliography}{\phantomsection\addcontentsline{toc}{section}{References}}
\setcounter{secnumdepth}{-1}
\setcounter{tocdepth}{2}
\allowdisplaybreaks[2]
\setlength{\emergencystretch}{3em}
\tolerance=2000
\begin{document}
\begin{center}
\vspace*{12pt}
{\LARGE Critical Dependence Stress for\\[5pt]
Diversification Gains under Fixed Marginals\par}
\vspace{18pt}
{\large Letao Ouyang\par}
\vspace{6pt}
{\small Independent Researcher, Nashua, NH 03062, USA\par}
\vspace{5pt}
{\small \href{mailto:theodore.oy2025@gmail.com}{\textcolor{black}{theodore.oy2025@gmail.com}}
\quad \href{mailto:10@alumni.duke.edu}{\textcolor{black}{10@alumni.duke.edu}}\par}
\vspace{10pt}
{\small Author manuscript\par}
\end{center}
\vspace{8pt}
\begin{abstract}
How much dependence uncertainty can a diversification benefit withstand when every marginal loss law is fixed? We study the smallest weighted $1$-Wasserstein perturbation that reverses the Expected Shortfall ranking of two nearby allocations evaluated under a common stressed law. For a positive continuous density on a compact box, exchangeability of the two traded positions eliminates the first-order gain, so the relevant stress radius is $\rho=s\varepsilon^2$ for trade size $\varepsilon$. At this scale, we establish a $\Gamma$-limit for the rescaled transport energies and prove convergence of the paired risk differences, with the two ES thresholds optimized separately. Coordinatewise quantile restoration incurs a marginal-flux cost, although its leading effects on the two ES values cancel in their difference. The limiting risk profile is continuous and strictly increasing, and its unique zero identifies the asymptotic critical budget. We call the excess over the local-efficiency lower bound the \emph{marginal-restoration premium}. This premium vanishes exactly when background displacement is cheapest per unit change in portfolio loss. Otherwise it is strictly positive, even when the profile is exactly linear on an initial interval. At equal effective costs, continuous maximum-flow/minimum-cut duality characterizes the right derivative of the critical budget with respect to the background penalty.
\end{abstract}
\begin{quote}\small
\noindent\textbf{Keywords:} Expected Shortfall; dependence uncertainty; full marginal constraints; Wasserstein distance; critical stress budget

\smallskip\noindent\textbf{MSC codes:} 91G70, 91G10, 49Q22.
\end{quote}
\newpage
\tableofcontents
\clearpage
\section{1. Introduction}\label{dependence-risk-in-a-reallocation-with-unchanged-acquisition-cost}

How much of a portfolio's diversification benefit survives uncertainty about dependence? A risk desk may have agreed marginal models for its individual positions without an equally well-determined model for their joint losses with the background book \cite{pflugpohl2018}. We ask how a proposed reallocation performs when these marginal models are retained but the coupling changes. Specifically, we seek the smallest dependence perturbation that reverses its Expected Shortfall (ES) advantage when both allocations face the same stressed law.

Let \(X_1,X_2\) be the unit losses of two positions of equal initial cost. A transfer from holdings \((1+\varepsilon,1-\varepsilon)\) to \((1,1)\), with an unchanged background loss \(\beta X_3\), changes the portfolio loss from \(A_\varepsilon=(1+\varepsilon)X_1+(1-\varepsilon)X_2+\beta X_3=B+\varepsilon H\) to \(B\), where
\[
B=X_1+X_2+\beta X_3,\qquad H=X_1-X_2,\qquad \beta>0.
\]
At upper-tail probability \(p\in(0,1)\), we compare them through \(\Phi_\varepsilon(P)=ES_p^P(B)-ES_p^P(A_\varepsilon)\), optimizing the two ES thresholds separately under \(P\) \cite{rockafellar2002}, so that \(\Phi_\varepsilon(P)<0\) favors the transfer under that law. Although \(\beta X_3\) cancels from the payoff difference \(A_\varepsilon-B\), its dependence with the traded positions determines which states enter each portfolio's upper tail. Put differently, a trade assessed in isolation can rank differently once embedded in the whole book \cite{brandtner2018}.

Writing \(\Pi\) for the Fr\'echet class of joint laws with these marginals, we measure distance from the reference law \(P_0\) by weighted \(W_1\) with ground cost \(c_\gamma(x,y)=|y-x|_\gamma=\sum_i\gamma_i|y_i-x_i|\), where each \(\gamma_i>0\) specifies the budget penalty per unit coordinate loss displacement as losses are re-paired across joint states. At budget \(\rho\), the stress problem is
\[
D_\varepsilon(\rho)
=\sup_{\substack{P\in\Pi\\ W_\gamma(P,P_0)\le\rho}}
\bigl\{ES_p^P(B)-ES_p^P(A_\varepsilon)\bigr\},
\qquad m_\varepsilon=-D_\varepsilon(0).
\]
Under an exchangeable reference law, the traded positions have equal marginal ES contributions, so the nominal gain is of second order: \(m_\varepsilon=(J/p)\varepsilon^2+o(\varepsilon^2)\), with \(J>0\). For example, when \(\gamma_1\le\min\{\gamma_2,\gamma_3/\beta\}\), the local exchange in Proposition~\ref{prop:nominal-benchmark} attains an initial stress derivative \(K_\varepsilon\) of order one, suggesting the reversal budget \(r_{0,\varepsilon}=m_\varepsilon/K_\varepsilon=O(\varepsilon^2)\). The derivative at zero budget does not determine whether this efficiency persists until the nominal gain is exhausted. Since the predicted reversal already occurs at scale \(\rho=s\varepsilon^2\), we need the risk profile over a nonvanishing interval of the rescaled budget \(s\), rather than its tangent at zero alone.

Three effects contribute at the critical scale. Moving the \(O(\varepsilon)\) mass in an \(O(\varepsilon)\)-layer around the reference tail boundary by distances of order \(\varepsilon\) changes both transport cost and tail expectation by \(O(\varepsilon^2)\), matching the nominal gain. This local movement need not preserve the marginals. For a fixed regular local plan, the marginal CDF discrepancies are \(O(\varepsilon^2)\); by Vallender's formula \cite{vallender1974}, their \(L^1\) norms give the one-dimensional restoration costs, even when compensation occurs far from the tail boundary. An \(O(\varepsilon)\) shift of either ES threshold also contributes at order \(\varepsilon^2\) through the curvature of its objective. We must therefore optimize the local displacement, marginal restoration, and both thresholds together.

When the reference law has a positive continuous density on a compact box \(I^2\times I_3\) and is exchangeable in its first two coordinates, Theorem~\ref{thm:critical-limit} gives a finite, continuous profile
\[
F(s)=\lim_{\varepsilon\downarrow0}\varepsilon^{-2}D_\varepsilon(s\varepsilon^2),
\qquad F(0)=-J/p,
\]
with convergence locally uniform in \(s\ge0\). The transport energies have a \(\Gamma\)-limit \cite{dalmaso1993} that includes both the local displacement cost and the cost of restoring its marginal projections. In the recovery argument, coordinatewise quantile maps attain the latter cost for each fixed regular local plan. Although this correction moves losses throughout the box, its leading effects on the two ES values cancel outside the threshold layer. Approximation then extends the paired-risk convergence to kernels with finite displacement moment.

Let \(\rho^-_\varepsilon\) and \(\rho^+_\varepsilon\) be the infima of budgets permitting a tie and a strict reversal. The strict-growth theorem identifies their common leading coefficient as the unique zero \(s^*\) of \(F\):
\[
\rho^-_\varepsilon=s^*\varepsilon^2+o(\varepsilon^2),\qquad
\rho^+_\varepsilon=s^*\varepsilon^2+o(\varepsilon^2).
\]
\emph{Initial efficiency and its limit.} The effective unit costs are \(\gamma_1,\gamma_2\), and \(c_3=\gamma_3/\beta\), with \(c_{\min}=\min\{\gamma_1,\gamma_2,c_3\}\). If the first traded coordinate is cheapest, Proposition~\ref{prop:balanced-linear-profile} constructs a positive interval \([0,s_{\rm lin}]\) on which \(F(s)=-J/p+s/(p\gamma_1)\). To reach zero at this rate, however, efficient displacement must remain possible throughout the required interval. Theorem~\ref{thm:zero-premium} characterizes attainment of the universal efficiency bound:
\[
s^*=Jc_{\min}\quad\Longleftrightarrow\quad
c_3\le\min\{\gamma_1,\gamma_2\}.
\]
We call the nonnegative gap \(s^*-Jc_{\min}\) the \emph{marginal-restoration premium} (normalized by \(Jc_{\min}\) when reported as a percentage). When background displacement is cheapest, exchangeability permits opposite shifts to balance on every fixed-background slice. When only traded directions are cheapest, a corner region requires a displacement moment of order \(\delta^2\), but its coordinate projections admit only \(O(\delta^4)\) of compensating moment near the diagonal (Figure~\ref{fig:corner}). Under the attained-slope condition of Proposition~\ref{prop:nominal-benchmark}, with \(c_{\min}=\gamma_1<c_3\), this gives \(\rho^\pm_\varepsilon/r_{0,\varepsilon}\to s^*/(\gamma_1J)>1\).

\emph{An analytical benchmark.} Consider independent uniform losses on \([0,2]^3\), with \(\beta=1\), \(p=1/2\), and \(\gamma=(1,1,2)\). Here \(g=3/8\) and \(J=5/32\), so Proposition~\ref{prop:nominal-benchmark} gives \(m_\varepsilon=(5/16)\varepsilon^2+o(\varepsilon^2)\). By Corollary~\ref{prop:uniform-variable-cost}, the profile is exactly \(F(s)=2s-5/16\) for \(0\le s\le1/96\), yet its zero satisfies \(13/64\le s^*\le29/96\). Even in this uniform model, the proved lower bound on reversal is already \(30\%\) above the local-efficiency prediction \(5/32\), despite an attainable initial slope.

\emph{Sensitivity at equal effective costs.} Along the cost family \(\gamma_1=\gamma_2=\gamma\), \(c_3=\gamma(1+\eta)\) for \(\eta\ge0\), the baseline \(Jc_{\min}=\gamma J\) is constant, so \(s^*\) and the marginal-restoration premium have the same right derivative \(\gamma Ja_+\) at \(\eta=0\). Here \(a_+\in(0,1)\) is the least background share among efficient stresses, characterized in Theorem~\ref{thm:projected-flow} by continuous maximum-flow/minimum-cut duality; Section 5 shows how this share and the budget sensitivity \(\gamma Ja_+\) can move in opposite directions as the reference dependence changes.

Our question lies between extremal risk aggregation and local distributional sensitivity. Subadditivity bounds aggregate risk \cite{artzner1999,mcneil2015}, while rearrangements within a Fr\'echet class describe how diversification depends on the coupling \cite{ruschendorf2013,embrechts2013,bernard2014,embrechts2015,wang2015superadditive}. Partial information leads to convex-order and mixability constraints \cite{li2018partial}, but can still permit comonotonic upper tails \cite{devecchinendelstreicher2026}, and fixed-marginal MES bounds quantify uncertainty in individual tail contributions \cite{chenfurmanlin2026}. For allocation comparisons, the symmetry imposed on the candidate laws is decisive. Chen et al.\ \cite{chen2022} and Liu and Liu \cite{liu2026} use exchangeability and majorization to identify diversification protection. Liu and Liu also analyze when complete dependence ambiguity favors concentration and how Wasserstein uncertainty modifies that conclusion. In our model, stressed couplings need not remain exchangeable. Symmetry at the reference gives stationarity, which does not by itself protect the reallocation once symmetry is broken under stress. We use this setting to measure the distance to a specified decision reversal, in the spirit of reverse sensitivity analysis \cite{pesenti2019,pesenti2022}.

Transport duality provides a natural way to preserve the marginal information \cite{villani2003,santambrogio2015}. In Wasserstein robust optimization \cite{blanchet2019,gao2023}, adding marginal constraints confines the perturbation to dependence changes \cite{gao2017,eckstein2020}. Such constraints encode calibrated risk-neutral laws in multiasset robust pricing \cite{eckstein2021multiasset}; Wasserstein and copula restrictions also arise in benchmark portfolio optimization \cite{pesentijaimungal2023}. Here the constraint acts on a common stressed loss law for two allocations. After rescaling, its full coordinate projections remain as a transport penalty, and the efficient allocation of displacement reduces to a capacity-constrained flow problem \cite{korman2015}.

Local Wasserstein expansions relate the sensitivity of optimized decisions and AVaR criteria to changes in the reference law \cite{bartl2021}, with adapted transport incorporating the information available in multiperiod decisions \cite{bartlwiesel2023}. When a first-order term vanishes, higher-order expansions become relevant. Fiechtner and Blanchet \cite{fiechtner2026} derive a second-order regret expansion under a unique empirical-risk minimizer in this situation, and Sauldubois \cite{sauldubois2026} develops higher-order Wasserstein and adapted-Wasserstein expansions, including martingale constraints. Dependence restrictions also affect local sensitivity, as in Lam \cite{lam2018}, while Namkoong \cite{namkoong2026} evaluates the first-order response over restricted transport directions. In our problem the trade size and stress radius vanish together, so locating a decision reversal at scale \(\rho=s\varepsilon^2\) requires the full rescaled profile, with global marginal compensation and both optimized ES thresholds retained over the budget interval needed to reach zero.

Section 2 establishes the nominal gain and the initial stress slope. Section 3 derives the variational limit at scale \(\rho=s\varepsilon^2\) and identifies the unique zero of the limiting profile. Section 4 characterizes zero premium and represents the boundary derivative by continuous projected flow--cut duality. Section 5 develops analytical examples and their implications for reverse stress testing. Appendix A contains the analytic proofs, and Appendix B gives the potential estimates and section calculations for the examples.

\paragraph{Notation and conventions} We write \(T_\#\mu\) for a pushforward and measure displacements by \(|d|_\gamma=\sum_i\gamma_i|d_i|\). The constant in \(O_\alpha(\varepsilon^k)\) may depend on the fixed kernel and model, but not on \(\varepsilon\). Scalar section integrands are extended by zero outside their domains, and projection identities are equalities of finite signed Borel measures. We use \(\inf\varnothing=+\infty\); allocation ratios are assigned the value zero where both numerator and denominator vanish.

\section{2. Model formulation and the local benchmark}\label{the-model-and-ranking-thresholds}

We first formulate the comparison for arbitrary integrable marginal laws, which suffice for existence of a worst stressed model. The critical-scale analysis will use a positive continuous density and exchangeability of the two traded positions, introduced in Section 2.2.

\subsection{2.1 Fixed marginals and a common stress law}
Let \(n\ge2\), let \(\mu_1,\ldots,\mu_n\) have finite first moments, and fix \(P_0\) in the Fr\'echet class \(\Pi(\mu_1,\ldots,\mu_n)=\{P\in\mathcal P(\mathbb R^n):P_i=\mu_i,\ 1\le i\le n\}\).

For exposure vectors \(b,h\in\mathbb R^n\), the losses before and after the proposed reallocation are
\(A_\varepsilon=B+\varepsilon H\) and \(B\), with \(B=b\cdot X\) and \(H=h\cdot X\).
We use the upper-tail mass \(p\in(0,1)\) and the Expected Shortfall representation of Rockafellar and Uryasev \cite{rockafellar2002},
\begin{equation}
ES_p^P(L)=\min_{t\in\mathbb R}
\left\{t+\frac1pE_P(L-t)_+\right\}.
\tag{2.1}\label{eq:legacy-2.1}
\end{equation}
We write \(\mathscr E_L(P;t)=t+p^{-1}E_P(L-t)_+\) for the threshold objective. For an integrable loss, its minimum is attained on the nonempty compact interval characterized by
\begin{equation}
t\in\arg\min_{u}\mathscr E_L(P;u)
\quad\Longleftrightarrow\quad P(L>t)\le p\le P(L\ge t).
\tag{2.2}\label{eq:legacy-A.1}
\end{equation}
Without a continuity assumption on the stressed loss law, support gaps can leave a nondegenerate interval of minimizing thresholds; we retain the set-valued characterization \eqref{eq:legacy-A.1} so that subsequent estimates apply to every minimizer, including for laws with atoms.

Under a common candidate law \(P\in\Pi\), we compare \(\Phi_\varepsilon(P)=ES_p^P(B)-ES_p^P(A_\varepsilon)\), minimizing each loss over its own threshold. Maximizing this paired difference tests whether a single dependence scenario can erase the nominal gain; separate maximizations of the two ES values may instead select different joint laws. We measure changes in dependence with positive weights \(\gamma_i\) and ground cost
\[
c_\gamma(x,y)=\sum_{i=1}^n\gamma_i|y_i-x_i|,
\qquad W_\gamma(P,P_0)=\min_{\pi\in\Gamma(P_0,P)}\int c_\gamma\,d\pi.
\]
Within \(\Pi=\Pi(\mu_1,\ldots,\mu_n)\), transport changes how component losses are paired across joint states. Holding the weights and all marginal laws fixed, we define the worst comparison at budget \(\rho\) and the nominal gain by

\begin{equation}
\mathcal A_\rho=\{P\in\Pi:W_\gamma(P,P_0)\le\rho\},\quad
D_\varepsilon(\rho)=\max_{P\in\mathcal A_\rho}\Phi_\varepsilon(P),\quad
m_\varepsilon=-\Phi_\varepsilon(P_0).
\tag{2.3}\label{eq:legacy-2.2}
\end{equation}
Proposition~\ref{prop:general-wellposedness} guarantees attainment in \eqref{eq:legacy-2.2}, allowing atoms, singular dependence, and unbounded marginal supports. Since a nondecreasing profile can vanish on an interval, we distinguish the infimal budgets for a tie and a strict reversal,
\begin{equation}
\rho^-_\varepsilon=\inf\{\rho\ge0:D_\varepsilon(\rho)\ge0\},
\qquad
\rho^+_\varepsilon=\inf\{\rho\ge0:D_\varepsilon(\rho)>0\},
\tag{2.4}\label{eq:legacy-2.3}
\end{equation}
which have the same leading asymptotics in the exchangeable box model of Section~2.2 by Theorem~\ref{thm:strict-crossing}.

\resultnumber{2.1}\begin{proposition}[Existence under integrable marginal laws]\label{prop:general-wellposedness}
Let \(n\ge2\), let \(\mu_1,\ldots,\mu_n\) have finite first moments, and fix \(P_0\in\Pi(\mu_1,\ldots,\mu_n)\), positive weights \(\gamma_i\), exposure vectors \(b,h\in\mathbb R^n\), and \(p\in(0,1)\). Then \(\Pi\) and every \(\mathcal A_\rho\) are compact for \(W_\gamma\), and the maximum in \eqref{eq:legacy-2.2} is attained. For each fixed \(\varepsilon\), the profile \(D_\varepsilon\) is nondecreasing and Lipschitz on \([0,\infty)\).
\end{proposition}
\begin{proof}
The fixed-marginal class is tight and weakly closed. Moreover, along any weakly convergent sequence \(P_k\to P\) in \(\Pi\), the first moments satisfy \(\int|x|_\gamma\,dP_k=\sum_i\gamma_i\int|x_i|\,d\mu_i=\int|x|_\gamma\,dP\), so the Wasserstein convergence criterion \cite{villani2003} implies \(W_\gamma\)-compactness of \(\Pi\) and its closed subsets \(\mathcal A_\rho\). For a linear loss \(L=a\cdot X\), coupling the laws and minimizing the uniformly Lipschitz hinges over their thresholds gives
\[
|ES_p^P(L)-ES_p^Q(L)|\le p^{-1}\|a\|_{\gamma,*}W_\gamma(P,Q),
\qquad \|a\|_{\gamma,*}=\max_i|a_i|/\gamma_i.
\]
Applying this bound to \(B=b\cdot X\) and \(A_\varepsilon=(b+\varepsilon h)\cdot X\), we find that their ES difference \(\Phi_\varepsilon\) is Lipschitz with constant \(C_\varepsilon=p^{-1}(\|b\|_{\gamma,*}+\|b+\varepsilon h\|_{\gamma,*})\). Compactness of \(\mathcal A_\rho\) then ensures attainment. To compare two budgets \(\rho_2>\rho_1\), take a maximizer \(P\) at \(\rho_2\). The mixture \(P_a=(1-a)P_0+aP\), \(a=\rho_1/\rho_2\), satisfies \(W_\gamma(P_a,P_0)\le a\rho_2\) and \(W_\gamma(P_a,P)\le(1-a)\rho_2\), by mixing an optimal plan with the corresponding identity plans. Comparing \(\Phi_\varepsilon(P_a)\) and \(\Phi_\varepsilon(P)\) yields \(0\le D_\varepsilon(\rho_2)-D_\varepsilon(\rho_1)\le C_\varepsilon(\rho_2-\rho_1)\).
\end{proof}

For a transfer \(H=X_i-X_j\) between positions with the same marginal law, replacing \(X_j\) by \(X_i\) preserves every marginal and gives \(A_\varepsilon=B\) at cost \(\gamma_jE_{P_0}|X_i-X_j|\); reversing their roles yields
\[
0\le\rho^-_\varepsilon
\le\min\{\gamma_i,\gamma_j\}\,E_{P_0}|X_i-X_j|<\infty.
\]
This guarantees a finite tie budget but does not identify its small-trade scale: under the exchangeable box assumptions introduced next, Section~3 shows that both critical budgets are of order \(\varepsilon^2\).

\subsection{2.2 The exchangeable box model}\label{model:critical-box}
To isolate how dependence with the rest of a portfolio affects a small reallocation, we keep the two traded losses separate and collect all unchanged positions into their aggregate loss \(\beta X_3\). We take the marginal law of \(X_3\) as retained information and specify the transport metric directly on \((X_1,X_2,X_3)\). This aggregate model fixes the information used in the stress test; its marginal and displacement cost are not equivalent to separately fixing the marginals and transport costs of every underlying background position.

\paragraph{Standing assumptions (exchangeable box model)}\label{assumption:critical-box}
Let
\[
I=[\ell,r],\qquad I_3=[\ell_3,r_3],\qquad
S=I^2\times I_3,\qquad \Delta_I=r-\ell>0,\quad \Delta_3=r_3-\ell_3>0.
\]
We also write \(I_1=I_2=I\) when treating coordinate projections together.
The reference law has a continuous density \(f>0\) on the closed box \(S\), normalized to integrate to one and satisfying
\(f(x,y,z)=f(y,x,z)\).
We take
\[
B=X_1+X_2+\beta X_3,\qquad H=X_1-X_2,\qquad
A_\varepsilon=B+\varepsilon H,\qquad
\beta>0,\quad 0<\varepsilon<1,
\]
with arbitrary \(p\in(0,1)\) and \(\gamma_i>0\), so that \(b=(1,1,\beta)\) and \(h=(1,-1,0)\) in the preceding notation. All model parameters remain fixed as \(\varepsilon\downarrow0\). The first two unit positions are expressed at equal initial cost, allowing the trade to change their holdings from \((1+\varepsilon,1-\varepsilon)\) to \((1,1)\) without changing the initial outlay. The position \(\beta X_3\) represents the background book, whose holding remains unchanged during the trade.

The stressed law may be any coupling in \(\Pi((P_0)_1,(P_0)_2,(P_0)_3)\); exchangeability and absolute continuity are assumptions on the reference law only. Reference symmetry makes the equal allocation stationary under \(P_0\), while equality of the first two marginals gives \(E_P[A_\varepsilon]-E_P[B]=\varepsilon(E_P[X_1]-E_P[X_2])=0\) under every candidate. Before implementation costs, the ES comparison therefore distinguishes tail risks between portfolios with the same expected loss.

The compact box and positive continuous density let us study this dependence mechanism with uniform control of marginal restoration. Positivity on the closed box gives both the interior mass needed for local exchanges and lower bounds on the marginal densities, hence Lipschitz bounds for their inverse CDFs. These bounds enter the quantile recovery in Section 3.2 and Appendix A.5; Section 5.3 considers the additional difficulties posed by unbounded marginals.

\paragraph{Loss units and effective stress costs}\label{m-units}
A shift \(\Delta X_3\) changes the portfolio loss by \(\beta\Delta X_3\) at cost \(\gamma_3|\Delta X_3|=(\gamma_3/\beta)|\beta\Delta X_3|\). Accordingly, \(c_1=\gamma_1\), \(c_2=\gamma_2\), and \(c_3=\gamma_3/\beta\) are the costs per unit change of \(B\) along the three coordinate axes. The ratio \(c_3/\gamma_1\) compares a background displacement with a first-position displacement having the same effect on portfolio loss.

\paragraph{Canonical box normalization} Writing the original coordinates as \(X=T_{\rm box}(u)=(\ell+\kappa u_1,\ell+\kappa u_2,\ell_3+\kappa_3u_3)\), with \(\kappa=\Delta_I/2\) and \(\kappa_3=\Delta_3/2\), gives a normalized model on \([0,2]^3\) with \(\widehat\beta=\beta\kappa_3/\kappa\) and \(\widehat\gamma=(\kappa\gamma_1,\kappa\gamma_2,\kappa_3\gamma_3)\). By cash invariance and positive homogeneity of ES, the original paired risk is \(\kappa\) times its normalized counterpart, while corresponding transport plans have equal cost; hence the critical budgets agree and \(p\) is unchanged. A common rescaling of all transport weights rescales those budgets.

\subsection{2.3 Nominal curvature and the exact initial slope}\label{nominal-improvement-initial-slope-and-the-protection-benchmark}

Under the reference model of Section 2.2, let \(v\) be the upper-tail \(p\) threshold of \(B\). The marginal expected shortfall (MES) of a unit position \(X_i\), namely \(E_{P_0}[X_i\mid B>v]\), is the ES derivative with respect to its holding, and multiplying by that holding gives its Euler contribution \cite{scaillet2004,tasche2001}. Since the first two holdings in \(B\) both equal one, the directional derivative along \(H=X_1-X_2\) is their MES difference: with \(T_H=E_{P_0}[H1_{\{B>v\}}]\), it equals \(p^{-1}T_H=E_{P_0}[X_1-X_2\mid B>v]=0\) by exchangeability. To determine the gain from a nonzero trade, we must then differentiate the moving tail boundary, which leads to a conditional variance on the tail section.

We record the classical section differentiation formula under our continuity assumptions in Lemma~\ref{lemma:general-variation}.

The tail boundary can be parameterized by
\[
\Omega=\{(x,y)\in(\ell,r)^2:\ell_3<z_v(x,y)<r_3\},\qquad
z_v(x,y)=\frac{v-x-y}{\beta},\qquad
\xi=(x,y,z_v(x,y)).
\]
\leavevmode\phantomsection\label{geometry:interior-diagonal}
Since \(p\in(0,1)\) and \(f>0\), the threshold satisfies \(2\ell+\beta\ell_3<v<2r+\beta r_3\). To locate a nondegenerate part of the section, follow the segment joining the lower and upper vertices of \(S\). The loss increases strictly along it and crosses \(v\) at an interior point \((x_*,x_*,z_*)\). A neighborhood of \((x_*,x_*)\) consequently lies in \(\Omega\), and its points off the diagonal have \(H=x-y\ne0\).

By the coarea formula, integration in the normal coordinate \(B-v\) gives the section measure
\(d\sigma=f(\xi)\,dx\,dy/\beta\). We denote the density of \(B\) at \(v\) and the nominal gain coefficient by
\begin{equation}
g=\sigma(\Omega)>0,\qquad J=\frac12\int_\Omega H^2\,d\sigma>0.
\tag{2.5}\label{eq:legacy-3.1}
\end{equation}
Here \(g=f_B(v)\) is the loss density at the threshold, and \(d\sigma/g\) is its normalized conditional section law. Exchangeability makes the conditional mean of \(H\) zero, giving \(2J/g=\operatorname{Var}_{P_0}(H\mid B=v)\). Phrased differently, the second-order diversification gain is governed by the conditional variance of the trade payoff on \(\{B=v\}\), weighted by the threshold density \(g\).

\resultnumber{2.2}\begin{proposition}[Nominal benchmark and exact initial efficiency]\label{prop:nominal-benchmark}
In the exchangeable box model of Section 2.2:
\begin{resultparts}
\item \emph{Nominal gain and threshold curvature.} The reference thresholds and nominal gain satisfy
\begin{equation}
t_\varepsilon-v=o(\varepsilon),\qquad
m_\varepsilon=\frac Jp\varepsilon^2+o(\varepsilon^2).
\tag{2.6}\label{eq:legacy-3.2}
\end{equation}
The threshold of \(B+\eta H\) is unique for every \(|\eta|<1\), and \(m_\varepsilon>0\) for \(0<\varepsilon<1\). With \(L_{c,\varepsilon}=B+c\varepsilon H\), the unoptimized objective satisfies, uniformly for \(c=0,1\) and \(|\tau|\le T<\infty\),
\begin{equation}
\mathscr E_{L_{c,\varepsilon}}(P_0;v+\varepsilon\tau)
=ES_p^{P_0}(B)+\frac{\varepsilon^2}{p}
\left(c^2J+\frac g2\tau^2\right)+o(\varepsilon^2).
\tag{2.7}\label{eq:legacy-B.2}
\end{equation}
\item \emph{Fixed-marginal Lipschitz bound and initial slope.} For every \(P,Q\in\Pi\),
\begin{equation}
|\Phi_\varepsilon(P)-\Phi_\varepsilon(Q)|\le \kappa_\varepsilon W_\gamma(P,Q),\quad
\kappa_\varepsilon=\frac1p\max\left\{\frac{1+\varepsilon/2}{\gamma_1},
\frac{1-\varepsilon/2}{\gamma_2},\frac\beta{\gamma_3}\right\}.
\tag{2.8}\label{eq:legacy-3.3}
\end{equation}
For fixed \(0<\varepsilon<1\), if
\begin{equation}
\frac{1+\varepsilon/2}{\gamma_1}
\ge\max\left\{\frac{1-\varepsilon/2}{\gamma_2},\frac\beta{\gamma_3}\right\},
\tag{2.9}\label{eq:legacy-3.4}
\end{equation}
then there exists \(r_\varepsilon>0\) such that
\begin{equation}
K_\varepsilon=\frac{1+\varepsilon/2}{p\gamma_1},\qquad
D_\varepsilon(\rho)=-m_\varepsilon+K_\varepsilon\rho
\quad(0\le\rho\le r_\varepsilon).
\tag{2.10}\label{eq:legacy-3.5}
\end{equation}
\end{resultparts}
\end{proposition}
\begin{proof}
For part~(i), the coefficients of \(B+\eta H\) are positive whenever \(|\eta|<1\), and its support is \([2\ell+\beta\ell_3,2r+\beta r_3]\). Since every interior level meets \(\operatorname{int}S\), Lemma~\ref{lemma:general-variation} gives a positive continuous loss density and a unique threshold. The gradients of \(H\) and \(B+\eta H\) are linearly independent, so \(H\) is nonconstant on each such section. Its conditional variance is therefore strictly positive, and the same lemma implies strict convexity of \(\eta\mapsto ES_p^{P_0}(B+\eta H)\) on \((-1,1)\).

At \(\eta=0\), exchanging the first two coordinates preserves both the reference tail and its boundary measure and reverses \(H\), giving \(T_H=m_1=0\). In the notation of Lemma~\ref{lemma:general-variation}, the remaining section moments at \((0,v)\) are \(m_0=g\) and \(m_2=2J\). Thus the threshold derivative vanishes and the optimized second derivative is \(2J/p\), which yield \eqref{eq:legacy-3.2} by Taylor expansion at zero. Strict convexity makes this stationary point the strict minimum, so \(m_\varepsilon>0\) for all \(0<\varepsilon<1\). For the unoptimized objective \(\Psi\), the Hessian at \((0,v)\) is \(\operatorname{diag}(2J/p,g/p)\). Expanding at the trade size and threshold displacement \((c\varepsilon,\varepsilon\tau)\) gives \eqref{eq:legacy-B.2}, uniformly on bounded \(\tau\)-intervals.

\emph{The fixed-marginal Lipschitz bound.} We next establish the uniform Lipschitz bound over \(\Pi\). Given any minimizing \(u\) for \(B\) under \(Q\) and any minimizing \(t\) for \(A_\varepsilon\) under \(P\), let
\[
\begin{aligned}
h_{u,t}&=(B-u)_+-(A_\varepsilon-t)_+,\\
\Delta_B&=\mathscr E_B(P;u)-ES_p^P(B),&
\Delta_A&=\mathscr E_{A_\varepsilon}(Q;t)-ES_p^Q(A_\varepsilon).
\end{aligned}
\]
Both gaps are nonnegative, and direct subtraction gives
\begin{equation}
\Phi_\varepsilon(P)-\Phi_\varepsilon(Q)
=p^{-1}\int h_{u,t}\,d(P-Q)-\Delta_B-\Delta_A.
\tag{2.11}\label{eq:legacy-A.14}
\end{equation}
A direct bound on \(h_{u,t}\) would use its largest absolute \(x_1\)-slope, \(1+\varepsilon\). Because \(P\) and \(Q\) share their marginals, adding an integrable separable potential does not change the integral against \(P-Q\). We can thus choose a representative of \(h_{u,t}\) modulo marginal potentials with a smaller Lipschitz bound. Here the required adjustment is linear: since \(\int H\,d(P-Q)=0\), we may add any scalar multiple of \(H=X_1-X_2\). The extreme \(x_1\)-slopes, \(-(1+\varepsilon)\) and \(1\), have midpoint \(-\varepsilon/2\), so adding \(\varepsilon H/2\) minimizes their largest absolute value among constant shifts and simultaneously centers the \(x_2\)-slopes. Away from the two tail planes, the four regime slopes are
\[
\begin{array}{c|rrrr}
(1_{\{B>u\}},1_{\{A_\varepsilon>t\}})&(0,0)&(1,0)&(0,1)&(1,1)\\
\partial_{x_1}h_{u,t}&0&1&-(1+\varepsilon)&-\varepsilon\\
\partial_{x_1}(h_{u,t}+\varepsilon H/2)&\varepsilon/2&1+\varepsilon/2&-1-\varepsilon/2&-\varepsilon/2.
\end{array}
\]
With the extreme \(x_1\)-slopes centered at zero, their largest absolute value drops to \(1+\varepsilon/2\), while the corresponding bounds for \(x_2\) and \(x_3\) are \(1-\varepsilon/2\) and \(\beta\). Setting \(G_{u,t}=h_{u,t}+\varepsilon H/2\), which is absolutely continuous along coordinate lines, we obtain
\[
|G_{u,t}(Y)-G_{u,t}(X)|
\le(1+\varepsilon/2)|Y_1-X_1|
+(1-\varepsilon/2)|Y_2-X_2|+\beta|Y_3-X_3|.
\]
Integrating against an optimal transport plan and then interchanging \(P,Q\) proves \eqref{eq:legacy-3.3}. The mixing argument in Proposition~\ref{prop:general-wellposedness}, now with this sharp constant, also gives
\begin{equation}
0\le D_\varepsilon(\rho_2)-D_\varepsilon(\rho_1)
\le \kappa_\varepsilon(\rho_2-\rho_1)\qquad(0\le\rho_1<\rho_2).
\tag{2.12}\label{eq:sharp-budget-continuity}
\end{equation}

\label{a.3-an-interior-intersection-and-attainment-of-the-initial-slope}
\emph{Attaining the initial slope.} Under \eqref{eq:legacy-3.4}, we construct an attaining exchange near an interior intersection of the reference tail planes \(\{B=v\}\) and \(\{A_\varepsilon=t_\varepsilon\}\). To locate such an intersection, fix \(\varepsilon\in(0,1)\) and let \([m_F,M_F]\) be the nondegenerate range of \(A_\varepsilon\) on \(F_v=\{B=v\}\cap S\). Both losses increase along each nonconstant segment from the lower vertex \(x_{\min}=(\ell,\ell,\ell_3)\). Any point with \(B>v\) lies beyond an intersection of that segment with \(F_v\), so \(t_\varepsilon\le m_F\) would imply \(\{B>v\}\subset\{A_\varepsilon>t_\varepsilon\}\). Since \(A_\varepsilon\) is nonconstant on \(F_v\), the second tail would also contain an open set below \(B=v\), contradicting the equal tail probabilities and \(f>0\). Rays toward the upper corner similarly exclude \(t_\varepsilon\ge M_F\), which leaves \(t_\varepsilon\in(m_F,M_F)\). The relative-interior property of a linear map on a convex section gives an intersection \(c=(x_*,y_*,z_*)\in\operatorname{int}S\).

To exchange the first-coordinate marginal intervals of two small source cubes in opposite disagreement regions along \(e_1\), we first choose reference points with the same coordinate \(x_*\) and seek equal and opposite changes in the two losses. From the interior intersection \(c\), a displacement \((0,-\zeta_0,a\zeta_0)\) changes \(B\) by \((-1+\beta a)\zeta_0\) and \(A_\varepsilon\) by \((-(1-\varepsilon)+\beta a)\zeta_0\). Equating their sum to zero determines \(a=(1-\varepsilon/2)/\beta\). Thus the points
\[
c_\pm=(x_*,y_*\mp \zeta_0,z_*\pm(1-\varepsilon/2)\zeta_0/\beta)
\]
have opposite strict tail memberships for small \(\zeta_0>0\): at \(c_+\), the losses differ from their thresholds by \(-\varepsilon \zeta_0/2\) for \(B\) and \(\varepsilon \zeta_0/2\) for \(A_\varepsilon\), with both signs reversed at \(c_-\). Choosing \(0<\delta_x<\varepsilon\zeta_0/(1+\varepsilon)\) bounds the loss change along either half-translation by the strict tail margin. With this margin fixed, we can take \(\eta_0>0\) small enough that the cubes \(E_\pm\) of half-width \(\eta_0\), centered at \(c_\pm\pm(\delta_x/2)e_1\), have disjoint second-coordinate intervals and remain inside \(S\) and their respective tail regions throughout translation by \(\mp\delta_xe_1\).

The uniform subdensities \(\nu_\pm=\lambda1_{E_\pm}\,dx\,dy\,dz\), where \(\lambda=\min_Sf\), have equal mass \(\mu=\lambda(2\eta_0)^3\) and sum at most \(P_0\). The translations \(T_\pm(x)=x\mp\delta_xe_1\) therefore define the admissible perturbation

\begin{equation}
P_\theta=P_0-\theta(\nu_++\nu_-)
+\theta(T_+)_\#\nu_++\theta(T_-)_\#\nu_-,\qquad 0\le\theta\le1.
\tag{2.13}\label{eq:legacy-A.16a}
\end{equation}
The translations exchange the first-coordinate marginal intervals, fix the other coordinates, and cross neither tail plane. Thus \(P_\theta\) preserves all marginals and both reference tail probabilities, leaving \(v\) and \(t_\varepsilon\) optimal. Integrating the affine hinges on the two cubes gives

\begin{equation}
\Phi_\varepsilon(P_\theta)-\Phi_\varepsilon(P_0)
=\frac{(2+\varepsilon)\theta\mu\delta_x}{p}
=K_\varepsilon(2\gamma_1\theta\mu\delta_x).
\tag{2.14}\label{eq:legacy-A.16b}
\end{equation}
The plan has cost \(2\gamma_1\theta\mu\delta_x\), and the sharp bound \eqref{eq:legacy-3.3} forces this to equal \(W_\gamma(P_\theta,P_0)\). Varying \(\theta\in[0,1]\) proves \eqref{eq:legacy-3.5} with \(r_\varepsilon=2\gamma_1\mu\delta_x>0\).
\end{proof}

If \(\gamma_2\ge\gamma_1\) and \(\gamma_3\ge\beta\gamma_1\), condition \eqref{eq:legacy-3.4} holds for every \(\varepsilon\in(0,1)\), suggesting the extrapolated reversal budget \(r_{0,\varepsilon}=m_\varepsilon/K_\varepsilon\). Because the disagreement regions narrow with \(\varepsilon\), this construction gives \(r_\varepsilon>0\) for each fixed trade size without controlling \(r_\varepsilon/\varepsilon^2\). Indeed, the unscaled bound \(|D_\varepsilon(\rho)|\le\varepsilon\|H\|_\infty\) makes both iterated limits in \((\rho,\varepsilon)\) vanish even though \(K_\varepsilon\to1/(p\gamma_1)>0\); Section 3 therefore tracks the profile under the joint scaling \(\rho=s\varepsilon^2\).

\paragraph{Reference symmetry and stressed dependence}\label{prop:exchangeable-protection}\label{rem:reference-symmetry}
Our stress class allows asymmetric dependence between the two traded positions. If every candidate law were exchangeable, the losses \(B-\varepsilon H\) and \(B+\varepsilon H\) would have the same distribution, and convexity of ES would give
\begin{equation}
ES_p^P(B)\le\tfrac12 ES_p^P(B-\varepsilon H)+\tfrac12 ES_p^P(B+\varepsilon H)
=ES_p^P(A_\varepsilon).
\tag{2.15}\label{eq:legacy-3.6}
\end{equation}
Requiring exchangeability of every stressed law would thus protect this reallocation \cite{chen2022,liu2026}. By imposing it only at the reference, we obtain stationarity there while permitting dependence stresses that test the persistence of the gain. Identical marginal laws alone do not imply stationarity. For example, on \(S=[0,2]^3\) with \(\beta=1\) and \(p=1/2\), the perturbation \(f_\delta=[1+\delta(x-y)(z-1)]/8\), with \(0<|\delta|<1/2\), preserves uniform marginals but gives a nonzero tail expectation \(T_H\). Antisymmetry in \(x,y\) preserves the law of \(B=x+y+z\), hence \(v=3\). Pairing \(z=1\pm t\) leaves the strip \(2-t<x+y\le2+t\), so
\[
T_H=\frac{\delta}{8}\int_0^1 t\int_{\{2-t<x+y\le2+t\}\cap[0,2]^2}(x-y)^2\,dx\,dy\,dt\ne0.
\]

\section{3. Variational limit and the critical profile}\label{critical-cost-and-recovery-under-full-marginal-constraints}

Under the traded-coordinate condition of Proposition~\ref{prop:nominal-benchmark}, a local exchange attains the initial slope. We now hold \(s=\rho/\varepsilon^2\) fixed to determine how far that efficiency persists toward reversal.
The limiting objects introduced below connect the stress profile to its reversal coefficient:
\[
\varepsilon^{-2}D_\varepsilon(s\varepsilon^2)
\xrightarrow{\text{locally uniformly}}
F(s)
\xrightarrow{\text{unique zero}}
s^*.
\]
The first passage follows from Theorem~\ref{thm:critical-limit}, which combines convergence of the transport energies with exact marginal recovery and convergence of the paired risk. Strict growth in Theorem~\ref{thm:strict-crossing} then identifies the common critical coefficient.

\subsection{3.1 Local transport and marginal compensation}\label{sec:limit-geometry}
An \(O(\varepsilon)\)-layer around \(\{B=v\}\) contains \(O(\varepsilon)\) reference mass, so displacements of order \(\varepsilon\) change both the transport cost and the paired hinge expectation at order \(\varepsilon^2\). We resolve this layer through the normal coordinate \(u=(B(X)-v)/\varepsilon\) and displacement \(d=(Y-X)/\varepsilon\), while retaining the cost of any transport outside bounded windows needed to restore the marginals.

\paragraph{Definition (Marginal flux and effective energy)} On bounded \(u\)-windows the rescaled source converges to \(d\nu=d\sigma\,du\). We disintegrate a limiting plan with this source into a probability kernel \(\alpha_{\xi,u}(dd)\), specifying the rescaled displacement at surface point \(\xi\) and normal coordinate \(u\); stationary sources have kernel \(\delta_0\). Writing \(d\Pi_\alpha(\xi,u,d)=\alpha_{\xi,u}(dd)\,d\nu(\xi,u)\), we require a finite weighted displacement moment even though the source measure \(\nu=\sigma\otimes du\) has infinite total mass:
\[
I(\alpha)=\int|d|_\gamma\,d\Pi_\alpha<\infty,
\qquad |d|_\gamma=\sum_i\gamma_i|d_i|.
\]
Across a coordinate cut \(\{X_i\le q\}\), upward movements remove mass from below the cut and downward movements add mass to it. The resulting signed flux has a Lebesgue density \(j_i\in L^1(I_i)\), as shown by the section formulas in Appendix A.3, and we define
\begin{equation}
j_i(q)dq=(\xi_i)_\#(d_i\,d\Pi_\alpha),\qquad
\mathcal C(\alpha)=I(\alpha)+\sum_{i=1}^3\gamma_i\|j_i\|_{L^1(I_i)}.
\tag{3.1}\label{eq:legacy-4.1}
\end{equation}
In words, \(j_i\) is the signed displacement moment projected onto the \(i\)-th coordinate, and its negative is the leading marginal CDF discrepancy produced by a regular local plan. Because the full plan preserves each marginal, transport outside a bounded normal and displacement window must compensate the discrepancy within it; the one-dimensional CDF formula for \(W_1\) \cite{vallender1974} gives the resulting cost lower bound. The energy \(\mathcal C\) includes this limiting restoration cost in addition to \(I\) (Figure~\ref{fig:tail-layer}). Projection contracts total variation, giving \(I\le\mathcal C\le2I\), with \(\mathcal C=I\) when all projected fluxes vanish and \(\mathcal C=2I\) when each \(d_i\) has a fixed sign \(\Pi_\alpha\)-almost everywhere, so that \(\|j_i\|_1=\int|d_i|\,d\Pi_\alpha\).

\begin{figure}[tbp]
\centering
\includegraphics[width=\linewidth]{figures/tail-layer-repair.pdf}
\caption{Local transport and full marginal compensation (schematic). A layer with thickness and mass of order \(O(\varepsilon)\), transported over distances of order \(O(\varepsilon)\), incurs second-order cost. The remaining nonnegative plan must cancel the local plan's coordinate CDF discrepancies. Their integrated absolute values give the compensation lower bound, with leading coefficients \(\|j_i\|_1\); a zero flux requires no corresponding leading-order compensation. The shaded discrepancy illustrates this mechanism.}
\label{fig:tail-layer}
\end{figure}


\subsection{3.2 The limiting risk and convergence theorem}\label{sec:limit-risk}
Under the same local plan, each loss has its own minimizing threshold. At threshold \(v+\varepsilon\tau\), a displacement \(d\) changes the leading centered and rescaled loss of \(L_{c,\varepsilon}=B+c\varepsilon H\) from \(u+cH\) to \(u+b\cdot d+cH\). Denote the corresponding hinge increment by \(\mathfrak q_c(\tau)=(u+b\cdot d+cH-\tau)_+-(u+cH-\tau)_+\). We write \(T_c^\alpha(\tau)=\int\mathfrak q_c(\tau)\,d\Pi_\alpha\) for its integral.

\paragraph{Definition (Limiting paired risk)} Moving the threshold away from \(v\) contributes the curvature term \(g\tau^2/2\) from \eqref{eq:legacy-B.2}. Adding the local transport increment gives the two limiting threshold objectives, which must still be minimized separately:
\begin{equation}
G_c^\alpha(\tau)=\frac g2\tau^2+T_c^\alpha(\tau),\qquad
\mathcal V(\alpha)=-\frac Jp+
\frac{\min_{\tau_B}G_0^\alpha(\tau_B)-\min_{\tau_A}G_1^\alpha(\tau_A)}p.
\tag{3.2}\label{eq:legacy-4.2}
\end{equation}

The energy \(\mathcal C\) depends only on the absolute displacement moment and the signed coordinate projections, and is convex under mixtures of kernels. The paired risk also depends on how the displacements are distributed relative to the normal coordinate \(u\): each of the two hinge objectives must optimize its own threshold, so their difference need not be concave in the common kernel.

We omit the superscript \(\alpha\) when the kernel is clear, and write \(\kappa_B=\max_i b_i/\gamma_i\), \(\kappa_H=\max\{\gamma_1^{-1},\gamma_2^{-1}\}\). Since \(|T_c^\alpha|\le\kappa_BI\), the quadratic term gives coercivity and confines every minimizer to \(g\tau^2/2\le2\kappa_BI\). Together with convexity from Lemma~\ref{lem:threshold-convexity}, these properties justify the threshold minimizations in \(\mathcal V\). We then optimize the paired value over the energy budget, with attainment established in Theorem~\ref{thm:critical-limit}:

\begin{equation}
F(s)=\max_{\mathcal C(\alpha)\le s}\mathcal V(\alpha),\qquad s\ge0.
\tag{3.3}\label{eq:legacy-4.3}
\end{equation}

We use vague convergence on \(\mathscr X=\mathcal M^+_{\rm loc}(Z)\), testing against compactly supported continuous functions on \(Z=I^2\times\mathbb R\times\mathbb R^3\), so that compensation may occur outside bounded windows. We lift a plan \(\pi\) from \(P_0\) to a law in \(\Pi\) by
\[
\mathscr T_\varepsilon\pi=\varepsilon^{-1}
\left(X_1,X_2,\frac{B(X)-v}{\varepsilon},\frac{Y-X}{\varepsilon}\right)_\#\pi.
\]
Here the factor \(\varepsilon^{-1}\) compensates for the thickness of the source layer. Because the lifted coordinates determine \((X,Y)\) for fixed \(\varepsilon\), we can retain the cost of the entire original plan in its rescaled energy:

\[
\begin{aligned}
\mathcal E_\varepsilon(\Lambda)
&=\begin{cases}
\displaystyle\int |d|_\gamma\,d\Lambda
=\varepsilon^{-2}\int c_\gamma\,d\pi,
&\begin{gathered}\Lambda=\mathscr T_\varepsilon\pi,\\ \pi\in\Gamma(P_0,P),\ P\in\Pi,\end{gathered}\\
+\infty,&\text{otherwise},
\end{cases}\\
\mathcal E(\Lambda)
&=\begin{cases}
\mathcal C(\alpha),&\Lambda=\Pi_\alpha,\ I(\alpha)<\infty,\\
+\infty,&\text{otherwise}.
\end{cases}
\end{aligned}
\]
Here we identify \(\Pi_\alpha\) with a measure on \(Z\) through \((\xi,u,d)\mapsto(\xi_1,\xi_2,u,d)\), using \(\xi=(x,y,z_v(x,y))\) on \(\Omega\) and assigning zero measure to \((I^2\setminus\Omega)\times\mathbb R\times\mathbb R^3\). Because \(\mathcal E_\varepsilon\) includes the whole plan, including transport outside bounded normal windows, minimizing it over plans with target law \(P\) gives \(\varepsilon^{-2}W_\gamma(P,P_0)\). We formulate the following \(\Gamma\)-convergence statement in this vague topology \cite{dalmaso1993}.

\keepproofstart{7\baselineskip}
\resultnumber{3.1}\begin{theorem}[\(\Gamma\)-convergence and the critical risk profile]\label{thm:critical-limit} Consider the exchangeable box model of \mbox{\hyperref[model:critical-box]{Section~2.2}}.
\begin{resultparts}
\item \emph{Transport plans and compactness.} \(\mathcal E_\varepsilon\) \(\Gamma\)-converges to \(\mathcal E\) as \(\varepsilon\downarrow0\) in the vague topology on \(\mathscr X\). Every sequence \(\varepsilon_n\downarrow0\), \(\Pi_n\in\mathscr X\), with bounded \(\mathcal E_{\varepsilon_n}(\Pi_n)\) has a vaguely convergent subsequence. If \(\Pi_n=\mathscr T_{\varepsilon_n}\pi_n\to\Pi_\alpha\) has bounded energy and \(P_n\) is the target law of \(\pi_n\), then \(\Phi_{\varepsilon_n}(P_n)/\varepsilon_n^2\to\mathcal V(\alpha)\).

\item \emph{Target laws and exact recovery.} For every sequence \(\varepsilon_n\downarrow0\) and every feasible sequence satisfying \(W_\gamma(P_n,P_0)=O(\varepsilon_n^2)\), there are a subsequence and a kernel with finite \(I\) such that

\begin{equation}
\mathcal C(\alpha)\le\liminf_n\frac{W_\gamma(P_n,P_0)}{\varepsilon_n^2},\qquad
\frac{\Phi_{\varepsilon_n}(P_n)}{\varepsilon_n^2}\longrightarrow\mathcal V(\alpha).
\tag{3.4}\label{eq:legacy-4.5}
\end{equation}

Conversely, every kernel with finite \(I\) admits a recovery sequence of distributions preserving the original three full marginals, with

\begin{equation}
\limsup_{\varepsilon\downarrow0}\frac{W_\gamma(P_\varepsilon,P_0)}{\varepsilon^2}\le\mathcal C(\alpha),\qquad
\frac{\Phi_\varepsilon(P_\varepsilon)}{\varepsilon^2}\longrightarrow\mathcal V(\alpha).
\tag{3.5}\label{eq:legacy-4.6}
\end{equation}

\item \emph{The stress profile.} The maximum in \eqref{eq:legacy-4.3} is attained. The function \(F\) is finite, continuous, and nondecreasing, with \(F(0)=-J/p\) and \(F>0\) at some finite budget. For every finite \(S_0\),

\begin{equation}
\sup_{0\le s\le S_0}\left|\frac{D_\varepsilon(s\varepsilon^2)}{\varepsilon^2}-F(s)\right|\longrightarrow0.
\tag{3.6}\label{eq:legacy-4.4}
\end{equation}
\end{resultparts}
\end{theorem}

\paragraph{Proof strategy}
Theorem~\ref{thm:critical-limit}(i) gives \(\Gamma\)-convergence of the transport energy and convergence of the paired risk along bounded-energy plans; exact marginal recovery in part~(ii) makes the limiting kernel values asymptotically attainable in \(\Pi\). Part~(iii) gives uniform convergence of the optimized profiles on compact budget intervals, which together with the strict crossing in Theorem~\ref{thm:strict-crossing} locates the reversal budget \(s^*\).

\emph{Step 1: Localization and the energy lower bound.} Because the two ES thresholds are optimized independently, comparison of the objectives requires a common threshold window for stresses with bounded rescaled energy. In Appendix~\hyperref[proof:threshold-localization]{A.2}, Lemma~\ref{b-l2-all-optimal-thresholds-under-a-budget-of-order-varepsilon2} places every minimizing threshold in an \(O_M(\varepsilon)\) window. On bounded windows the lifted plans are locally compact; outside them, marginal conservation requires compensation for the projected local flux. The compensation bound above then gives the energy lower bound for arbitrary admissible plans. Choosing optimal plans yields the corresponding Wasserstein bound for their target laws.

\emph{Step 2: Convergence of the paired risk.} Bounded first moments make each hinge increment integrable, but need not give uniform integrability along convergent plans. The two objectives share the same large-displacement contribution, so we subtract \(\mathfrak q_0(0)\) from both. The relative increments are bounded on the common threshold window. Outside a normal window enlarged to account for the bounded reallocation \(H\), a nonzero relative increment requires a sufficiently large displacement, so the first moment controls the far-field contributions. The baseline integral is finite but need not converge. Since it translates both threshold objectives by the same amount independently of the threshold, it cancels exactly from their difference of minima.

\emph{Step 3: Exact marginal recovery.} Let \(\mathcal K_{\rm fin}^c\) denote probability kernels with compactly supported continuous weights on finitely many fixed displacement vectors and the remaining mass at zero, as in \eqref{eq:legacy-B.26}. For a fixed kernel in this class, local translations yield a preliminary law \(Q_\varepsilon\) with uniform marginal CDF expansion \(F_{i,\varepsilon}-F_i=-\varepsilon^2j_i+o(\varepsilon^2)\), where \(F_i\) is the reference CDF. The quantile maps \(F_i^{-1}\circ F_{i,\varepsilon}\) restore all three marginals at cost \(\varepsilon^2\sum_i\gamma_i\|j_i\|_1+o(\varepsilon^2)\), with \(O_\alpha(\varepsilon^2)\) displacements.

To see why this correction does not alter the leading paired risk, compare its displacement scale with the \(O_\alpha(\varepsilon)\) threshold band. Outside a sufficiently wide band, the paired hinge has gradient \(-\varepsilon h\) above both thresholds and zero below both, so multiplying by the restoration displacement gives an \(O_\alpha(\varepsilon^3)\) contribution. Inside the band the gradient is bounded, but the \(Q_\varepsilon\)-mass is only \(O_\alpha(\varepsilon)\), giving the same order after integration. The restoration error is therefore \(O_\alpha(\varepsilon^3)\), uniformly over all minimizing threshold pairs.

Approximation by \(\mathcal K_{\rm fin}^c\) in cost, paired risk, and the vague plan topology (Appendix A.4), followed by diagonal recovery, gives the general \(o(\varepsilon^2)\) error. Compactness, contraction toward the identity, and the common budget-Lipschitz bound then give locally uniform convergence of the optimized profiles.


\subsection{3.3 Initial linearity and the critical budget}\label{sec:balanced-profile}
\phantomsection\label{g.4-balanced-source-selection-and-the-limiting-linear-segment}
At zero threshold shifts, the portfolios disagree on tail membership precisely on \(W=\{u(u+H)<0\}\). The displacement \(\zeta=-(u+H)1_W\) moves a source in this wedge to the limiting \(A_\varepsilon\)-boundary and attains the maximal paired hinge increment per unit of normal movement. When \(c_{\min}=\gamma_1\), we can assign movement along \(e_1\) provided its downward and upward moments balance on each \(x\)-fiber. On active sources in \(W\), \(\zeta<0\) for \(y<x\) and \(\zeta>0\) for \(y>x\), with \(\int\zeta\,du=-H|H|/2\). Writing \(k(x,y)=f(x,y,(v-x-y)/\beta)1_{\{\beta\ell_3<v-x-y<\beta r_3\}}\), the available fiber moments are
\begin{equation}
\ell_-(x)=\frac1{2\beta}\int_I(x-y)_+^2k(x,y)\,dy,\qquad
\ell_+(x)=\frac1{2\beta}\int_I(y-x)_+^2k(x,y)\,dy.
\tag{3.7}\label{eq:legacy-A.19}
\end{equation}
Although exchangeability gives \(\int_I\ell_-=\int_I\ell_+=J/2\), the two moments need not agree on each fiber. Matching their pointwise minimum cancels the projected flux, yielding the budget
\[
\begin{aligned}
q_{\rm bal}(x)&=\min\{\ell_-(x),\ell_+(x)\},&
Q_{\rm bal}&=\int_I q_{\rm bal}(x)\,dx,\\
s_{\rm lin}&=2\gamma_1Q_{\rm bal}=\gamma_1J(1-\chi)>0,&
\chi&=\|\ell_--\ell_+\|_{L^1(I)}/J.
\end{aligned}
\]
\resultnumber{3.2}\begin{proposition}[Capacity of the initial linear segment]\label{prop:balanced-linear-profile} In the exchangeable box model of Section~2.2, suppose \(\gamma_1=\min\{\gamma_1,\gamma_2,\gamma_3/\beta\}\). With \(\ell_\pm\) as in \eqref{eq:legacy-A.19} and \(s_{\rm lin}=\gamma_1J(1-\chi)>0\) as above,
\begin{equation*}
F(s)=-J/p+s/(p\gamma_1),\qquad 0\le s\le s_{\rm lin}.
\end{equation*}
\end{proposition}

\begin{proof}
For \(W=\{u(u+H)<0\}\) and \(\zeta=-(u+H)1_W\), select the proportion
\[
a(\xi)=
\begin{cases}
q_{\rm bal}(x)/\ell_-(x),&H>0,\ \ell_-(x)>0,\\
q_{\rm bal}(x)/\ell_+(x),&H<0,\ \ell_+(x)>0,\\
0,&\text{otherwise},
\end{cases}
\]
and define \(\alpha^\theta_{\xi,u}=\theta a(\xi)\delta_{\zeta e_1}+(1-\theta a(\xi))\delta_0\) for \(0\le\theta\le1\). Because \(0\le a\le1\) and the active sources satisfy \(|u|,|\zeta|\le\Delta_I\), this is a finite-first-moment probability kernel. Its opposite first-coordinate fluxes equal \(-\theta q_{\rm bal}(x)\) and \(\theta q_{\rm bal}(x)\), while the other coordinates remain fixed, so \(j_i=0\) for all \(i\) and \(I(\alpha^\theta)=\mathcal C(\alpha^\theta)=2\theta\gamma_1Q_{\rm bal}=\theta s_{\rm lin}\).

Under this displacement, an active source has target normal losses \(-H\) for \(B\) and zero for \(A_\varepsilon\). Because \(\operatorname{sgn}u=\operatorname{sgn}(-H)\ne0\) on \(W\), its strict \(B\)-tail membership is unchanged and no mass is placed at the \(B\)-threshold. For \(A_\varepsilon\), however, the active sources from the two signs of \(H\) arrive at the kink itself. The resulting atom accounts for the jump between the one-sided derivatives: with \(\mathfrak m_\pm=\int_{\{\pm H>0\}}a(\xi)|H|\,d\sigma\),
\[
G'_{0,-}(0)=G'_{0,+}(0)=0,\qquad
G'_{1,-}(0)=-\theta \mathfrak m_-\le0,\qquad
G'_{1,+}(0)=\theta \mathfrak m_+\ge0.
\]
By Lemma~\ref{lem:threshold-convexity}, zero is globally optimal for both objectives. Evaluating the nonzero wedge integrals, \(H^2/2\) for \(B\) on \(H<0\) and \(-H^2/2\) for \(A_\varepsilon\) on \(H>0\), we obtain \(T_0(0)=\theta Q_{\rm bal}\) and \(T_1(0)=-\theta Q_{\rm bal}\), which yield \(\mathcal V(\alpha^\theta)=-J/p+\theta s_{\rm lin}/(p\gamma_1)\). Since \(0<\chi<1\) by Lemma~\ref{lem:projected-imbalance}, we may choose \(\theta=s/s_{\rm lin}\) to attain the upper bound inherited from Proposition~\ref{prop:nominal-benchmark}.
\end{proof}

\paragraph{Strict growth} To identify a common leading budget for a tie and a strict reversal, we improve every bounded-first-moment kernel at arbitrarily small additional cost. Mixing with a kernel of larger paired risk would give a direct route if the paired risk were concave in the kernel. The two threshold objectives are minimized separately, however, and their difference need not be concave along mixtures. Nor can we assume that a kernel leaves stationary mass available for an added displacement: an everywhere nonzero shift proportional to \(e^{-u^2}e_1\) already has finite first moment.

We instead replace a small fraction of the existing kernel on two exchanged normal bands beyond the common threshold window. Reflecting \(u\) to \(-u\) crosses that window. Pairing the bands cancels their \(B\)-increments, while opposite signs of \(H\) produce a positive paired gain per unit replacement fraction, independent of the distance \(U\). After subtraction of the common baseline, the old kernel contributes only an \(O_M(U^{-1})\) remainder, because affecting the threshold window from those bands requires a large displacement. We first choose \(U\) so this remainder is smaller than the reflection gain, then reduce the replacement fraction to meet the additional budget. Appendix A.6 verifies the mixing obstruction and carries out this comparison uniformly over the original and modified minimizing thresholds.

\resultnumber{3.3}\begin{theorem}[Strict improvement by local replacement and a common critical coefficient]\label{thm:strict-crossing} Consider the exchangeable box model of \mbox{\hyperref[model:critical-box]{Section~2.2}}. The limiting profile \(F\) is strictly increasing and has a unique zero \(s^*>0\), which gives the common leading coefficient of both pre-limit critical budgets at scale \(\varepsilon^2\). More precisely:\begin{resultparts}\item \emph{Quantitative local improvement.} For every \(0\le M<\infty\) there exists \(c_M>0\) such that every kernel satisfying \(I(\alpha)\le M\) and every \(0<\delta\le1\) admit a feasible kernel \(\widehat\alpha\) with

\begin{equation}
\mathcal C(\widehat\alpha)\le\mathcal C(\alpha)+\delta,\qquad
\mathcal V(\widehat\alpha)\ge\mathcal V(\alpha)+c_M\delta.
\tag{3.8}\label{eq:legacy-5.1}
\end{equation}

These kernels can be chosen as a family \((\widehat\alpha_\delta)_{0<\delta\le1}\) satisfying
\[
\int |d|_\gamma\,|\widehat\alpha_\delta-\alpha|(dd)\,d\nu
\xrightarrow[\delta\downarrow0]{}0,
\qquad
\lim_{\delta\downarrow0}
\frac{\mathcal V(\widehat\alpha_\delta)-\mathcal V(\alpha)}{\delta}
\ge c_M>0.
\]
For every finite \(S_0>0\),

\begin{equation}
F(t)-F(s)\ge c_{S_0}(t-s)>0\qquad(0\le s<t\le S_0).
\tag{3.9}\label{eq:legacy-5.2}
\end{equation}

\item \emph{Critical budgets.} Together with the continuity, negative initial value, and positive value at a finite budget in Theorem~\ref{thm:critical-limit}(iii), strict growth of \(F\) gives a unique zero \(s^*>0\). Over the finite-first-moment kernel class,

\begin{equation}
s^*=\min_{\alpha:\,\mathcal V(\alpha)\ge0}\mathcal C(\alpha)
=\min_{\alpha:\,\mathcal V(\alpha)=0}\mathcal C(\alpha)
=\inf_{\alpha:\,\mathcal V(\alpha)>0}\mathcal C(\alpha).
\tag{3.10}\label{eq:legacy-5.3}
\end{equation}
The corresponding pre-limit critical budgets satisfy
\begin{equation}
\rho^-_\varepsilon=s^*\varepsilon^2+o(\varepsilon^2),\qquad
\rho^+_\varepsilon=s^*\varepsilon^2+o(\varepsilon^2).
\tag{3.11}\label{eq:legacy-5.4}
\end{equation}

The first two minima are attained, whereas the infimum over strictly positive risk values is not.
\end{resultparts}\end{theorem}

\paragraph{Execution costs}\label{cor:m-c1} A positive execution cost proportional to \(\varepsilon\) exceeds the second-order nominal gain for small trades. To retain a critical-budget comparison, suppose instead that the deterministic cost paid to implement the new allocation \(B\) satisfies \(c_{\mathrm{exec},\varepsilon}\ge0\) and \(c_{\mathrm{exec},\varepsilon}/\varepsilon^2\to\theta\in[0,J/p)\). With the stress class and metric fixed, cash invariance gives \(D^{\rm net}_\varepsilon=D_\varepsilon+c_{\mathrm{exec},\varepsilon}\), whose rescaled profile converges locally uniformly to \(F+\theta\).

Applying the tie and strict-reversal definitions in \eqref{eq:legacy-2.3} to the net profile, we obtain by strict growth and local uniform convergence \(\rho^{\pm,\rm net}_\varepsilon=s_\theta\varepsilon^2+o(\varepsilon^2)\), where \(s_\theta\) is the unique solution of \(F(s_\theta)=-\theta\). As the rescaled execution cost \(\theta\) increases to \(J/p\), this solution decreases continuously from \(s^*\) to zero. For example, in the cheap-background case of Section~4.1, the explicit linear profile gives \(s_\theta=c_3(J-p\theta)\).

\section{4. Efficiency premiums and projected flow duality}\label{background-substitutability-and-minimum-cost-feasibility}

We now turn to the question of when the critical coefficient reaches the local-efficiency lower bound and when a strict marginal-restoration premium remains. Because a locally efficient displacement need not preserve the marginal projections, any unbalanced moment must be compensated elsewhere. We measure this obstruction by testing each projected flux against a bounded measurable coordinate potential, using the \(L^1\) dual representation of the restoration cost:
\[
\mathcal C(\alpha)-I(\alpha)
=\sup_{|\phi_i|\le\gamma_i}\sum_i\int_{I_i}\phi_i(q)j_i(q)\,dq
=\sup_{|\phi_i|\le\gamma_i}\int\sum_i\phi_i(\xi_i)d_i\,d\Pi_\alpha.
\]
Here \(\phi_i(q)\) acts as a dual shadow price for the projected displacement moment at coordinate level \(q\). Vanishing projections require opposite signed displacements to balance at each coordinate level, eliminating the leading restoration penalty. We first identify the pointwise displacement required for equality in the efficiency bound, then test its projections near an off-diagonal corner to rule out equality when only traded coordinates are cheapest. At equal effective costs, continuous potentials yield the flow--cut duality of Section~4.2. Allowing nonzero projected fluxes in the restoration-cost bound leads to the piecewise-constant potential estimates of Section~5.1.

\subsection{4.1 Local efficiency and the zero-premium classification}\label{a-necessary-and-sufficient-condition-for-zero-premium}\label{v21-classification}

With \(c_1=\gamma_1\), \(c_2=\gamma_2\), \(c_3=\gamma_3/\beta\), and \(c_{\min}=\min\{c_1,c_2,c_3\}\), Proposition~\ref{prop:nominal-benchmark} and Theorem~\ref{thm:critical-limit} already give \(F(s)\le-J/p+s/(pc_{\min})\). The key issue at zero risk difference is whether a kernel can satisfy the pointwise efficiency conditions while conserving every marginal projection. We characterize those simultaneous equality conditions below.

Proposition~\ref{prop:nominal-benchmark} establishes initial efficiency at each fixed trade size under \eqref{eq:legacy-3.4}; when \(\gamma_1=c_{\min}\), Proposition~\ref{prop:balanced-linear-profile} separately realizes the limiting slope by matching opposite first-coordinate moments on each fiber. When only traded coordinates are cheapest, that matching cannot cover the whole tail section, although it suffices for the initial linear segment. If the background coordinate is cheapest, exchanging \(x,y\) and reversing \(u\) instead balances the required displacement on every background slice, sustaining the efficiency bound up to zero risk.

\resultnumber{4.1}\begin{lemma}[Local-efficiency bound and critical-level optimality]\label{lem:efficiency-bound}\label{lem:v21-equality}\label{g.3-a-necessary-and-sufficient-condition-for-minimum-cost-feasibility}
In the exchangeable box model of Section~2.2, every kernel with \(I(\alpha)<\infty\) satisfies \(p\mathcal V(\alpha)+J\le\mathcal C(\alpha)/c_{\min}\), so \(s^*\ge c_{\min}J\). Set \(b=(1,1,\beta)\) and \(\zeta=-(u+H)1_{\{u(u+H)<0\}}\). Then \(s^*=c_{\min}J\) if and only if there is a measurable probability kernel satisfying
\begin{equation}
\begin{aligned}
b\cdot d=\zeta,\qquad |d|_\gamma=c_{\min}|\zeta|
&\quad\Pi_\alpha\text{-almost everywhere},\\
(\xi_i)_\#(d_i\,d\Pi_\alpha)=0
&\quad i=1,2,3.
\end{aligned}
\tag{4.1}\label{eq:legacy-6.3}
\end{equation}
Such a kernel has \(\mathcal C=I=c_{\min}J\) and \(\mathcal V=0\); every kernel with \(\mathcal V\ge0\) and \(\mathcal C\le c_{\min}J\) satisfies these conditions.
\end{lemma}

\begin{proof}
Set \(w=b\cdot d\), \(\Delta_+(t,w)=(t+w)_+-t_+\), and choose any \(\tau_A\in\arg\min G_1\). To compare the risk gain with its transport cost, we evaluate the \(B\)-objective at zero, bound the difference of hinge increments by \(|w|\), and use \(c_{\min}|w|\le |d|_\gamma\). Recording the loss in each comparison gives
\begin{equation}
\begin{aligned}
\frac{\mathcal C}{c_{\min}}-(p\mathcal V+J)
={}&\underbrace{\frac{\mathcal C-I}{c_{\min}}}_{\mathrm{I}}
+\underbrace{\left(\frac I {c_{\min}}-\int |w|\,d\Pi_\alpha\right)}_{\mathrm{II}}
+\underbrace{\frac g2\tau_A^2}_{\mathrm{III}}\\
&+\underbrace{\int[|w|-\Delta_+(u,w)+\Delta_+(u+H-\tau_A,w)]\,d\Pi_\alpha}_{\mathrm{IV}}\\
&+\underbrace{G_0(0)-\min_{\tau_B}G_0(\tau_B)}_{\mathrm{V}}.
\end{aligned}
\tag{4.2}\label{eq:legacy-E.5}
\end{equation}
The five terms (I)--(V) in \eqref{eq:legacy-E.5} have distinct roles: (I) the marginal-restoration penalty; (II) the excess transport cost over the cheapest normal movement, including cancellation between coordinate displacements; (III) the curvature cost of shifting the \(A_\varepsilon\)-threshold; (IV) the unsaturated hinge gain; and (V) the gap from evaluating the \(B\)-objective at zero instead of its minimum. For the fourth term, \(\Delta_+(t,w)\in[\min\{0,w\},\max\{0,w\}]\), an interval of length \(|w|\), so the difference of the two hinge increments cannot exceed \(|w|\). All five terms are consequently nonnegative, proving the efficiency bound. Since \(\tau_A\) was an arbitrary minimizer, equality forces each defect to vanish for every optimal \(A_\varepsilon\)-threshold.


To reach zero risk at the local-efficiency budget, none of these defects can remain. Indeed, if \(\mathcal V\ge0\) and \(\mathcal C\le c_{\min}J\), the left side of \eqref{eq:legacy-E.5} is nonpositive, so every nonnegative term must vanish and
\begin{equation}
\mathcal V=0,\qquad \mathcal C=I=c_{\min}J,\qquad
\int|w|\,d\Pi_\alpha=J,\qquad j_i=0\quad(i=1,2,3).
\tag{4.3}\label{eq:legacy-E.7}
\end{equation}
In addition, \(\arg\min G_1=\{0\}\) and \(0\in\arg\min G_0\). The vanishing metric defect
\begin{equation}
|d|_\gamma-c_{\min}|b\cdot d|
=\sum_i(\gamma_i-c_{\min}b_i)|d_i|
+c_{\min}\left(\sum_i b_i|d_i|-\left|\sum_i b_id_i\right|\right)
\tag{4.4}\label{eq:legacy-E.8}
\end{equation}
allows only cheapest coordinates, with all nonzero components having the sign of \(w\). At \(\tau_A=0\), equality in the hinge bound requires \(u\ge0\), \(u+H\le-w\) if \(w>0\), and \(u\le0\), \(u+H\ge-w\) if \(w<0\). Thus \(|w|\le K:=|u+H|1_{\{u(u+H)<0\}}\), with the sign of \(\zeta\) wherever \(w\ne0\). Since
\begin{equation}
\int_{\mathbb R}K\,du=H^2/2,\qquad
\int_{\mathbb R}\zeta\,du=-H|H|/2,
\tag{4.5}\label{eq:legacy-E.12}
\end{equation}
and every conditional kernel has mass one, \(\int K\,d\Pi_\alpha=J=\int|w|\,d\Pi_\alpha\). Saturation of this nonnegative bound forces
\begin{equation}
b\cdot d=\zeta\quad\Pi_\alpha\text{-almost everywhere},
\tag{4.6}\label{eq:legacy-6.4}
\end{equation}
which, with the metric and flux equalities, proves necessity.

For the converse, the threshold calculation depends only on the normal displacement \(w=\zeta\). We integrate over the disagreement wedge and use exchangeability to obtain
\begin{equation}
T_0(0)=J/2,\qquad T_1(0)=-J/2,\qquad
\int_{H>0}H^2\,d\sigma=\int_{H<0}H^2\,d\sigma=J.
\tag{4.7}\label{eq:legacy-D.2}
\end{equation}
Because the transformed \(B\)-loss retains its strict tail sign and places no atom at zero, we have \(G_0'(0)=0\). With \(m_\pm=\int(\pm H)_+\,d\sigma>0\), the mass placed at the \(A_\varepsilon\)-threshold kink gives
\begin{equation}
G_{1,-}'(0)=-m_-<0,\qquad G_{1,+}'(0)=m_+>0.
\tag{4.8}\label{eq:legacy-D.3}
\end{equation}
By convexity, these strict one-sided inequalities make zero the unique minimizer of \(G_1\). For \(G_0\), the stationary regions \(\{H>0,u>0\}\) and \(\{H<0,u<0\}\) each have section mass \(g/2\), and their unchanged normal coordinates cover the positive and negative half-lines, respectively. Hence \(\lambda_0\ge(g/2)\,dt\) on \(\mathbb R\), and Lemma~\ref{lem:threshold-convexity} yields \(G_0(\tau)-G_0(0)\ge g\tau^2/4\), so zero is also the unique minimizer of \(G_0\). Combining these minima with \eqref{eq:legacy-D.2} gives \(\mathcal V=0\). Under \eqref{eq:legacy-6.3}, we have \(I=c_{\min}J\) and vanishing fluxes, so \(\mathcal C=c_{\min}J\). Together with attainment of the critical cost in Theorem~\ref{thm:strict-crossing}, this proves the equivalence.
\end{proof}

\resultnumber{4.2}\begin{theorem}[Zero-premium classification]\label{thm:zero-premium} For the exchangeable box model of \mbox{\hyperref[model:critical-box]{Section~2.2}} and every positive weight vector,
\begin{equation}
s^*=Jc_{\min}\quad\Longleftrightarrow\quad
c_3\le\min\{\gamma_1,\gamma_2\}.
\tag{4.9}\label{eq:legacy-6.1}
\end{equation}
If \(c_3=c_{\min}\), the profile is explicitly \(F(s)=-J/p+s/(pc_3)\) for \(0\le s\le c_3J\).

If \(\gamma_2\ge\gamma_1\) and \(c_3>\gamma_1\), then
\begin{equation}
1<\frac{s^*}{\gamma_1J}\le\min\{1+\chi,c_3/\gamma_1\}<2,
\tag{4.10}\label{eq:legacy-6.2}
\end{equation}
where \(\chi=\|\ell_--\ell_+\|_{L^1(I)}/J\in(0,1)\) is the projected imbalance of Section~3.3.
\end{theorem}

\keepproofstart{4\baselineskip}
\textit{Proof of Theorem~\ref{thm:zero-premium}.}

By Lemma~\ref{lem:efficiency-bound}, it suffices to allocate the prescribed displacement \(\zeta\) among cheapest coordinates while preserving every marginal projection. The background coordinate admits such a balance because exchanging the traded losses and reversing the normal coordinate leaves each background slice fixed and reverses the required displacement. For the traded coordinates, the corresponding balance must instead be found along each \(x\)- or \(y\)-fiber. Near an off-diagonal corner, the required moment is bounded below by a positive multiple of \(\delta^2\), whereas the opposite-side sources available at those coordinate levels have capacity only \(O(\delta^4)\). We prove the background construction first and then use this discrepancy to exclude even randomized splitting among the traded directions.

\paragraph{Comparison kernels and the sufficient direction}\label{d.1-three-zero-risk-difference-kernels-scores-and-optimal-thresholds}\label{d.2-net-fluxes-total-costs-and-general-bounds}
Assigning the prescribed displacement to one coordinate gives
\begin{equation}
\alpha^{(1)}=\delta_{\zeta e_1},\qquad
\alpha^{(2)}=\delta_{\zeta e_2},\qquad
\alpha^{(3)}=\delta_{(\zeta/\beta)e_3}
\tag{4.11}\label{eq:legacy-D.1}
\end{equation}
with finite first moments and zero risk differences by the normal-displacement calculation \eqref{eq:legacy-D.2}--\eqref{eq:legacy-D.3} in Lemma~\ref{lem:efficiency-bound}. The first two kernels have respective fluxes \(j_1=\ell_+-\ell_-\) and \(j_2=\ell_--\ell_+\), all other fluxes being zero. To verify the background projection, we use the exchange \((x,y,u)\mapsto(y,x,-u)\), which preserves \(\nu,\xi_3\) and reverses \(\zeta\). Absolute integrability then allows us to test against every bounded measurable \(\varphi\), obtaining
\begin{equation}
\int\varphi(\xi_3)d_3^{(3)}d\nu
=-\int\varphi(\xi_3)d_3^{(3)}d\nu=0.
\tag{4.12}\label{eq:legacy-D.6}
\end{equation}
All three projections of \(\alpha^{(3)}\) vanish, yielding
\begin{equation}
\mathcal C(\alpha^{(1)})=\gamma_1J(1+\chi),\quad
\mathcal C(\alpha^{(2)})=\gamma_2J(1+\chi),\quad
\mathcal C(\alpha^{(3)})=c_3J.
\tag{4.13}\label{eq:legacy-D.7}
\end{equation}
The background kernel proves sufficiency, including ties, and the least of these costs gives \(s^*\le J\min\{\gamma_1(1+\chi),\gamma_2(1+\chi),c_3\}\), which yields the upper bound in \eqref{eq:legacy-6.2} when \(\gamma_2\ge\gamma_1\) and \(c_3>\gamma_1\).

When \(c_3=c_{\min}\), mix the background kernel \(\alpha^{(3)}\) with the identity using weight \(\theta=s/(c_3J)\in[0,1]\). Both kernels have zero projected flux and both threshold objectives attain their minima at zero, by \eqref{eq:legacy-D.2}--\eqref{eq:legacy-D.3} and \(G_c^{\delta_0}(\tau)=g\tau^2/2\), so cost and paired risk are affine along the mixture. Its cost is \(\theta c_3J=s\), and \eqref{eq:legacy-D.2} gives
\begin{equation}
\mathcal V=-(1-\theta)J/p.
\tag{4.14}\label{eq:legacy-D.9}
\end{equation}
Thus the mixture attains the efficiency bound on the stated interval.

\paragraph{The obstruction for arbitrary randomized kernels}\label{e.3-a-local-obstruction-to-zero-flux-in-the-two-asset-directions}
For necessity, suppose \(c_3>c_{\min}\) and a kernel with \(\mathcal V\ge0\) attains cost \(c_{\min}J\). Lemma~\ref{lem:efficiency-bound} then requires \(d_3=0\), \(d_1+d_2=\zeta\), and both traded components to have the sign of \(\zeta\). Since the kernel has finite first moment, its averaged displacement fields \(m_i(x,y)=\int_{\mathbb R}\int d_i\,\alpha_{\xi,u}(dd)\,du\), \(i=1,2\), are integrable on the section and satisfy, by \eqref{eq:legacy-6.3},
\begin{equation}
\begin{aligned}
&m_1+m_2=-H|H|/2,\qquad |m_i|\le H^2/2,\\
&m_i\ge0\quad\text{for}\ x<y,\qquad
m_i\le0\quad\text{for}\ x>y.
\end{aligned}
\tag{4.15}\label{eq:legacy-E.15}
\end{equation}
Projected conservation gives, for every pair of Borel sets \(A,B\subset I\),
\begin{equation}
\int_{\Omega\cap\{x\in A\}}m_1d\sigma=0,\qquad
\int_{\Omega\cap\{y\in B\}}m_2d\sigma=0.
\tag{4.16}\label{eq:legacy-E.16}
\end{equation}
We test \eqref{eq:legacy-E.15}--\eqref{eq:legacy-E.16} near the upper-left off-diagonal corner of \(\Omega\), where \(x\) is small, \(y\) is large, and the positive moment \((y-x)^2/2\) is bounded away from zero. Compensating this moment at the same coordinate levels requires negative moment from \(\{x>y\}\), where the small value of \(x\) or large value of \(y\) confines available sources to a small triangle near the diagonal (Figure~\ref{fig:corner}).

\leavevmode\phantomsection\label{lem:v21-corner-geometry}
The section base is the intersection of an open square with a diagonal strip, \(\Omega=(\ell,r)^2\cap\{q_-<x+y<q_+\}\), whose boundaries and coordinate projection endpoints are
\[
\begin{gathered}
q_-=v-\beta r_3,\qquad q_+=v-\beta\ell_3,\\
\xi_-=\max\{\ell,q_--r\},\qquad
\xi_+=\min\{r,q_+-\ell\},\qquad \Delta=\xi_+-\xi_->0.
\end{gathered}
\]
Both coordinate projections of \(\Omega\) equal \((\xi_-,\xi_+)\), and the interior diagonal point from \mbox{\hyperref[geometry:interior-diagonal]{Section~2.3}} satisfies \(\xi_-<x_*<\xi_+\), confirming \(\Delta>0\). Moreover, \(q_*:=\xi_-+\xi_+\in[q_-,q_+]\), so the corner \((\xi_-,\xi_+)\) lies in the closed strip, possibly on its boundary. The square diagonal through that corner runs parallel to the strip boundaries. On the corner square \((x,y)=(\xi_-+t,\xi_+-z)\), \(0<t,z<\delta\), the strip condition reads \(-(q_*-q_-)<t-z<q_+-q_*\). Because the two nonnegative slacks sum to \(\beta\Delta_3\), at least one is at least \(\beta\Delta_3/2\). Choosing \(\delta<\beta\Delta_3/2\) therefore includes one of the triangles \(z<t\) and \(t<z\). If we also choose \(\delta<\Delta/4\), we obtain
\[
E_\delta=\Omega\cap((\xi_-,\xi_-+\delta)\times(\xi_+-\delta,\xi_+)),
\qquad |E_\delta|\ge\delta^2/2,
\qquad y-x\ge\Delta/2\quad\hbox{on }E_\delta.
\]
Because \(f\) is bounded below on \(S\) and \(y-x\ge\Delta/2\) on this region, there is a constant \(c_\Omega>0\), independent of \(\delta\), such that
\begin{equation}
M_\delta:=\int_{E_\delta}(m_1+m_2)d\sigma
\ge c_\Omega\delta^2.
\tag{4.17}\label{eq:legacy-E.20}
\end{equation}
Testing \eqref{eq:legacy-E.16} on \(A_\delta=(\xi_-,\xi_-+\delta)\) and \(B_\delta=(\xi_+-\delta,\xi_+)\) requires \(M_\delta\le N_1+N_2\), where
\[
N_1=\int_{\Omega\cap\{x\in A_\delta,\ x>y\}}(-m_1)d\sigma,
\qquad
N_2=\int_{\Omega\cap\{y\in B_\delta,\ x>y\}}(-m_2)d\sigma.
\]
In \(N_1\), the conditions \(x\in A_\delta\) and \(y<x\) force \(\xi_-<y<x<\xi_-+\delta\); in \(N_2\), the sources similarly satisfy \(\xi_+-\delta<y<x<\xi_+\). Each triangle has area \(O(\delta^2)\), and \eqref{eq:legacy-E.15} bounds the corresponding displacement moment by \(-m_i\le H^2/2=(x-y)^2/2=O(\delta^2)\). Applying the upper density bound gives
\begin{equation}
N_1+N_2\le C_\Omega\int_0^\delta\int_0^x(x-y)^2\,dy\,dx
=\frac{C_\Omega}{12}\delta^4.
\tag{4.18}\label{eq:legacy-E.22}
\end{equation}
Here \(C_\Omega<\infty\) is independent of \(\delta\). For small \(\delta\), the bound contradicts \(c_\Omega\delta^2\le M_\delta\le N_1+N_2\). Attainment of the critical cost then implies \(s^*>Jc_{\min}\) when \(c_3>c_{\min}\), and \eqref{eq:legacy-D.7} gives \eqref{eq:legacy-6.2}. \unskip\nobreak\hfill\proofbox

% V21_CORNER_FIGURE
\begin{figure}[tbp]
\centering
\includegraphics[width=\linewidth]{figures/corner-capacity.pdf}
\caption{Corner demand and return capacity, illustrated on \(I=I_3=[0,2]\) for \(\beta=7/4\), \(v=15/4\), \(\delta=0.35\). The blue region \(E_\delta\) stays away from the diagonal and has demand at least a positive multiple of \(\delta^2\). Return sources lie in the hatched triangles, which may extend beyond \(\Omega\); their area and squared displacement are both \(O(\delta^2)\), giving capacity \(O(\delta^4)\). The bounds shown, with \(\lambda=\min f\), \(\Lambda=\max f\), and \(\Delta=\xi_+-\xi_-\), violate \(M_\delta\le N_1+N_2\) for sufficiently small positive \(\delta\).}
\label{fig:corner}
\end{figure}


\subsection{4.2 Boundary sensitivity and the flow--cut representation}\label{required-background-share-and-the-boundary-response}\label{v21-flow}\label{variational-characterization-and-bounds-for-the-background-share}
At equal effective costs, every efficient kernel has total required displacement \(J\). At this critical level, the equality conditions and total-moment saturation in Lemma~\ref{lem:efficiency-bound} fix \(b\cdot d=\zeta\) pointwise almost everywhere, so normal integration and exchange symmetrization reduce the remaining minimum-background allocation problem to shares on the two-dimensional half-section. Increasing only the background penalty selects, to first order, those kernels using the least background displacement. We can calculate this amount by assigning as much displacement as possible to the traded coordinates under their projected conservation law. On \(D=\Omega\cap\{x<y\}\), the measure \(d\mu=w\,d\sigma=w_\sigma(x,y)\,dx\,dy\), where \(w=(y-x)^2/2\), has mass \(J/2\) and a continuous density, positive in the interior.

\resultnumber{4.3}\begin{theorem}[Boundary derivative and projected flow duality]\label{prop:boundary-sensitivity}\label{thm:projected-flow}
Fix the reference model and \(p,\beta,\gamma\), and set
\[
\boldsymbol\gamma(\eta)=(\gamma,\gamma,\beta\gamma(1+\eta)),\quad
\pi(\eta)=s^*(\boldsymbol\gamma(\eta))/(\gamma J)-1,\quad \eta\ge0.
\]
Let \(\mathcal C_0\) be the cost at \(\eta=0\), and let \(\mathcal Z_0=\{\alpha:I(\alpha)<\infty,\ \mathcal V(\alpha)=0,\ \mathcal C_0(\alpha)=\gamma J\}\). When kernels are identified with their measures \(\Pi_\alpha\), the set \(\mathcal Z_0\) is nonempty and compact in the vague topology on \(\mathscr X\) from Section~3.2. The following minimum is attained, and
\begin{equation}
\pi(\eta)=a_+\eta+o(\eta),\qquad
a_+=\frac\beta J\min_{\alpha\in\mathcal Z_0}
\int|d_3|\,d\Pi_\alpha\in[0,1].
\tag{4.19}\label{eq:legacy-6.5}
\end{equation}

On \(D=\Omega\cap\{x<y\}\), equip the half-section with \(d\mu=\frac12(y-x)^2\,d\sigma\). For measurable shares \(a,b\), define
\begin{equation*}
\mathfrak F_\mu=\max\left\{\int_D(a+b)d\mu:
a,b\ge0,\ a+b\le1,\ (x)_\#(a\mu)=(y)_\#(b\mu)\right\}.
\end{equation*}
The optimum is attained and satisfies \(0<\mathfrak F_\mu<J/2\). Moreover,
\begin{equation}
a_+=1-\frac{2\mathfrak F_\mu}{J}\in(0,1),\qquad
\mathfrak F_\mu=\inf_{\substack{\varphi\text{ measurable}\\|\varphi|\le1}}
\int_D\max\{0,1-\varphi(x),1+\varphi(y)\}\,d\mu.\tag{4.20}\label{eq:legacy-6.8}
\end{equation}
Equivalently, \(a_+\) is the supremum of the projected cut deficits
\begin{equation}
a_+=\frac2J\sup_{U\cap V=\varnothing}
\left[\mu(D\cap(U\times V))-
\mu\bigl(D\cap((V\times I)\cup(I\times U))\bigr)\right].\tag{4.21}\label{eq:legacy-6.9}
\end{equation}
Here \(U,V\subset I\) are disjoint measurable sets.

\end{theorem}

In words, \(2\mathfrak F_\mu\) is the largest normal-displacement moment that the traded coordinates can accommodate while balancing their marginal projections. The cut formula \eqref{eq:legacy-6.9} expresses \(a_+\) as the maximal deficit between the required moment on \(D\cap(U\times V)\) and the compensating capacity on \(D\cap((V\times I)\cup(I\times U))\), multiplied by \(2/J\).

\textit{Proof sketch.} For \(\alpha\in\mathcal Z_0\), normal integration and exchange symmetrization yield shares satisfying \((x)_\#(a\mu)=(y)_\#(b\mu)\) and background amount \(\beta\int|d_3|\,d\Pi_\alpha=J-2\int_D(a+b)\,d\mu\). Conversely, every feasible pair of shares lifts to a kernel in \(\mathcal Z_0\), so the minimum background amount is \(J-2\mathfrak F_\mu\) (Appendix~\hyperref[app:flow-proof]{A.8}). Dividing by \(J\) gives \(a_+\), and Appendix~\hyperref[app:cost-proofs]{A.7} establishes \eqref{eq:legacy-6.5} by comparing perturbed minimizers with a least-background kernel. Sion's theorem with continuous coordinate potentials gives \eqref{eq:legacy-6.8}; signed layer-cake reduction then gives \eqref{eq:legacy-6.9}.

\paragraph{Dependence on the transport weights}\label{d.4-weights-the-tail-surface-and-the-order-of-limits}
Varying the positive transport weights leaves the finite-first-moment kernel class and \(\mathcal V\) unchanged. With \(A_i(\alpha)=\int|d_i|\,d\Pi_\alpha+\|j_i\|_1\), the coordinate cost including marginal restoration, we may therefore write
\begin{equation}
s^*(\gamma)=\inf_{\mathcal V(\alpha)\ge0}\sum_i\gamma_i A_i(\alpha).
\tag{4.22}\label{eq:legacy-D.11}
\end{equation}
Because the feasible class is independent of \(\gamma\), this infimum of nonnegative linear functions is positively homogeneous, coordinatewise nondecreasing, and concave. In particular, \(r_-\gamma\le\widetilde\gamma\le r_+\gamma\), with \(0<r_-\le r_+\), implies \(r_-s^*(\gamma)\le s^*(\widetilde\gamma)\le r_+s^*(\gamma)\), ensuring stability under relative weight perturbations.

The strict-premium classification also persists as long as the perturbed background cost remains above the cheapest traded cost. For example, equal effective costs \(\gamma=(1,1,\beta)\) lie on the zero-premium boundary. Starting instead from \(\beta=1\) and \(\gamma=(1,1,2)\), any perturbation with \(r_+<2r_-\) satisfies \(\widetilde\gamma_3>\max\{\widetilde\gamma_1,\widetilde\gamma_2\}\) and therefore retains a strict premium.

\keepproofstart{6\baselineskip}
\section{5. Analytical illustrations and financial implications}\label{critical-cost-premia-background-dependence-and-finite-reallocations}\label{v21-quantitative}
\subsection{5.1 Projection potentials and a uniform reference}\label{strict-premium-despite-equal-efficiency}

By assigning a bounded shadow price \(\phi_i\) to each marginal flux, we can bound the strict premium without solving the full risk optimization. To derive the estimate, we first minimize the adjusted coordinate costs \(\gamma_i\pm\phi_i(\xi_i)\) for a prescribed normal displacement and then compare that minimum with the paired hinge gain.

\resultnumber{5.1}\begin{proposition}[A projection-potential bound]\label{prop:projection-potential-bound}
In the exchangeable box model, let \(b=(1,1,\beta)\) and let \(|\phi_i|\le\gamma_i\) be bounded measurable potentials on the coordinate intervals such that
\(k_\pm(\xi)=\min_{1\le i\le3}(\gamma_i\pm\phi_i(\xi_i))/b_i\)
are bounded below by a positive constant. For every finite-first-moment kernel, every \(\tau\in\arg\min G_1\), and every \(\lambda>0\),
\begin{equation}
\mathcal C(\alpha)\ge
\lambda\left(J+p\mathcal V+\frac g2\tau^2\right)
-\frac{\lambda}{2}\int_\Omega(H-\tau)^2
\left(\frac{\lambda}{k(\xi,H-\tau)}-1\right)_+\,d\sigma,
\tag{5.1}\label{eq:legacy-J.6}
\end{equation}
where \(k(\xi,h)=k_+(\xi)\) for \(h<0\) and \(k_-(\xi)\) for \(h>0\). The integrand is zero when \(h=0\).
\end{proposition}

The proof in \hyperref[app:potential-proof]{Appendix B.1} optimizes the hinge gain against the adjusted cost \(k_+w_++k_-(-w)_+\), where \(w=b\cdot d\). To quantify how far the reversal budget can exceed the initial-slope prediction, we use an analytically tractable reference with independent uniform losses on \([0,2]\), \(\beta=1\), \(p=1/2\), and \(\gamma=(\gamma_0,\gamma_0,2\gamma_0)\). Equal exposures and the median tail make the section invariant under central reflection, so \(v=3\), \(d\sigma=dx\,dy/8\), \(g=3/8\), and \(J=5/32\).

In the cut representation of Theorem~\ref{thm:projected-flow}, the midpoint partition corresponds to \(\varphi(q)=1_{[0,1)}(q)-1_{(1,2]}(q)\). Scaling this potential by \(\gamma_0/2\) and using opposite signs on the two traded coordinates transfers the same partition to the marginal-flux penalty in Proposition~\ref{prop:projection-potential-bound}, giving the admissible two-level potentials
\[
\phi_1(x)=\frac{\gamma_0}{2}\operatorname{sgn}(1-x),\qquad
\phi_2(y)=\frac{\gamma_0}{2}\operatorname{sgn}(y-1),\qquad \phi_3=0.
\]
Their signs distinguish low-loss from high-loss coordinate halves. At zero threshold shift, a positive paired hinge gain requires normal movement \(w=b\cdot d>0\) when \(x<y\), and \(w<0\) when \(x>y\). The amplitude \(\gamma_0/2\) gives effective cost \(k(\xi,H)=3\gamma_0/2\) in this direction across the midpoint and both directional costs \(\gamma_0/2\) on either same-side region. Choosing \(\lambda=3\gamma_0/2\) in \eqref{eq:legacy-J.6} therefore removes the cross-region penalty at zero threshold. By retaining the correction for shifted minimizing thresholds in the proof below, we obtain a bound valid without prescribing either ES threshold.

\resultnumber{5.2}\begin{corollary}[Uniform reference with costly background displacement]\label{prop:uniform-variable-cost}
For this uniform reference,
\begin{equation}
\frac{13}{64}\le\frac{s^*}{\gamma_0}\le\frac{29}{96},
\qquad \frac{s^*}{\gamma_0J}\ge\frac{13}{10}.
\tag{5.2}\label{eq:legacy-7.1}
\end{equation}
The initial linear profile holds on \(0\le s\le\gamma_0/96\).
\end{corollary}

\begin{proof}
For \(1/2\le x\le1\), the negative-moment section has \(1-x<y<x\), so integrating \((x-y)^2/16\) over this interval gives \(\ell_-(x)=(2x-1)^3/48\). For \(0\le x\le1/2\), the section is empty and \(\ell_-(x)=0\). Reflection gives \(\ell_+(x)=\ell_-(2-x)\). On the lower half of \(I\), \(\ell_-\) is the smaller moment, which gives \(Q_{\rm bal}=2\int_{1/2}^1\ell_-(x)\,dx=1/192\). Using \(\int_I(\ell_-+\ell_+)\,dx=J\) then gives \(\|\ell_--\ell_+\|_1=J-2Q_{\rm bal}=7/48\), and Proposition~\ref{prop:balanced-linear-profile} yields \(s_{\rm lin}=\gamma_0/96\). The first-coordinate kernel \(\delta_{\zeta e_1}\) from \eqref{eq:legacy-D.1} reaches zero risk at cost \(\gamma_0(J+\|\ell_--\ell_+\|_1)=29\gamma_0/96\). For the lower bound, the midpoint potentials with \(\lambda=3\gamma_0/2\) give
\[
\frac{\mathcal C}{\gamma_0}\ge
\begin{cases}
\displaystyle \frac34\mathcal V+\frac{13}{64}
+\frac{\tau^2(6-\tau^2)}{64},& |\tau|\le1,\\[2pt]
\displaystyle \frac12\mathcal V+\frac{11}{32},& |\tau|\ge1.
\end{cases}
\]
For \(|\tau|\le1\), the section integrations in \hyperref[app:uniform-bound-proof]{Appendix B.2} give the first line. For \(|\tau|\ge1\), the efficiency decomposition \eqref{eq:legacy-E.5} gives the second. Every kernel with \(\mathcal V\ge0\) therefore costs at least \(13\gamma_0/64\), which is \(30\%\) above \(\gamma_0J\).
\end{proof}

\phantomsection\label{app:quantitative-benchmarks}
Setting \(\gamma_0=1\), we obtain analytical envelopes for \(F(s)\) on \([0,U]\), where \(U=29/96\). The balanced kernel has cost \(s_0=1/96\), zero projected flux, and risk \(-7/24\), while the first-coordinate comparison kernel has cost \(U\) and risk zero. Their convex mixtures have affine paired risk because zero minimizes both threshold objectives for each kernel, and affine cost because the balanced kernel has zero flux. Joining these mixtures to the exact initial linear segment, we realize the lower envelope
\[
\underline F(s)=
\begin{cases}
2s-5/16,&0\le s\le1/96,\\
s-29/96,&1/96\le s\le29/96.
\end{cases}
\]
To bound \(F\) from above, we invert the two threshold estimates and combine them with the efficiency bound:
\[
\begin{aligned}
F(s)&\le\min\left\{2s-\frac5{16},
\max\left\{\frac43s-\frac{13}{48},2s-\frac{11}{16}\right\}\right\}
\\
&=\overline F(s):=\min\left\{2s-\frac5{16},\frac43s-\frac{13}{48}\right\}
\qquad(0\le s\le U).
\end{aligned}
\]
The last equality holds because the branches inside the maximum meet only at \(s=5/8>U\). Figure~\ref{fig:budget} plots the analytical envelopes \(\underline F\le F\le\overline F\) on \([0,U]\). The guaranteed linear segment \([0,s_{\rm lin}]\) covers one fifteenth of the tangent prediction \(J=5/32\), whereas the critical coefficient \(s^*\) lies between \(13/10\) and \(29/15\) times \(J\). Thus the initial slope is attained on a positive interval, yet the required reversal budget is at least \(30\%\) higher than its tangent prediction.
\begin{figure}[tbp]
\centering
\includegraphics[width=\linewidth]{figures/budget-comparison.pdf}
\caption{Proved risk envelopes for the independent uniform model with \(\beta=1\), \(p=1/2\), and \(\gamma=(1,1,2)\). Explicit mixtures attain the lower envelope; the upper envelope follows from midpoint potentials and the efficiency bound. The inset shows the exact initial interval \([0,1/96]\) in green. The dashed line extrapolates its slope to \(J=5/32\), whereas the critical budget lies in \([13/64,29/96]\). Both plotted envelopes are piecewise linear analytical formulas.}
\label{fig:budget}
\end{figure}


\subsection{5.2 A dependence family with fixed marginal laws}\label{background-dependence-changes-substitutability-even-with-identical-marginals}

Section~5.1 used an independent uniform reference with background displacement twice as costly as either traded direction. We now move to the equal-effective-cost boundary and ask how changing the reference dependence affects the required background share \(a_+\) and the budget derivative \(\gamma_0Ja_+\). Retaining uniform marginals, \(\beta=1\), and \(p=1/2\), we perturb the independent density by the product of the centered losses \(x+y-2\) and \(z-1\):
\[
f_\vartheta(x,y,z)=\frac{1+\vartheta(x+y-2)(z-1)}8,\qquad
|\vartheta|<\frac12,\qquad
E_\vartheta[(X_1+X_2-2)(X_3-1)]=\frac{2\vartheta}{9}.
\]
After rescaling the losses to \([0,1]\), this is a one-parameter subfamily of the trivariate Farlie--Gumbel--Morgenstern family \cite[Section~2]{otakimura2021}: the interactions between each traded coordinate and the background have parameter \(\vartheta\), with the other interaction parameters zero. Centering preserves every marginal, and the parameter restriction keeps the density strictly positive on the cube. Central reflection fixes \(v=3\), where the section density is \([1-\vartheta(x+y-2)^2]/8\), so increasing \(\vartheta\) reduces the section mass away from \(x+y=2\) even as covariance increases. On the unchanged half-section \(D\), let
\(d\mu_\vartheta=w_{\sigma,\vartheta}(x,y)\,dx\,dy\), where
\(w_{\sigma,\vartheta}(x,y)=(y-x)^2\{1-\vartheta(x+y-2)^2\}/16\).
Set \(D_-=D\cap\{y<1\}\), \(D_+=D\cap\{x>1\}\), and denote the same-side masses by \(M_{\vartheta,\pm}=\mu_\vartheta(D_\pm)\).

\resultnumber{5.3}\begin{proposition}[Background share along the dependence family]\label{prop:dependence-family}
For \(|\vartheta|<1/2\), the boundary flow is attained and
\begin{equation}
J_\vartheta=\frac{75-14\vartheta}{480},\qquad
M_{\vartheta,-}=M_{\vartheta,+}=\frac{3-2\vartheta}{1152},
\qquad
a_{+,\vartheta}
=1-\frac{4(M_{\vartheta,-}+M_{\vartheta,+})}{J_\vartheta}.
\tag{5.3}\label{eq:legacy-K.9}
\end{equation}
In particular,
\[
J_\vartheta a_{+,\vartheta}=\frac{195-22\vartheta}{1440},\qquad
\frac{d a_{+,\vartheta}}{d\vartheta}
=\frac{360}{(75-14\vartheta)^2}>0,\qquad
\frac{d(J_\vartheta a_{+,\vartheta})}{d\vartheta}
=-\frac{11}{720}.
\]
\end{proposition}

To obtain the explicit value, we use all the mass in \(D_-\) and \(D_+\) and match their projections in two disjoint parts of the cross-region. The section inequalities in \hyperref[app:dependence-family-proof]{Appendix B.3} provide the required capacity throughout \(|\vartheta|<1/2\). Because the midpoint potential gives the same value in the dual problem, this feasible flow is optimal.

As reference dependence varies through \(0<\vartheta<1/2\), the minimum background share and the corresponding unnormalized displacement moment change in opposite directions:
\begin{equation}
a_{+,\vartheta}-a_{+,0}
=\frac{24\vartheta}{5(75-14\vartheta)}>0,\qquad
J_\vartheta a_{+,\vartheta}-J_0a_{+,0}
=-\frac{11\vartheta}{720}<0.
\tag{5.4}\label{eq:legacy-7.2}
\end{equation}
The reason for this contrast is that, on the median plane \(x+y+z=3\), the centered factors satisfy \(z-1=-(x+y-2)\). Thus increasing \(\vartheta\) favors equal signs in the full box but reduces the section density away from \(x+y=2\) by the factor \(1-\vartheta(x+y-2)^2\). Across both exchanged half-sections, the traded-coordinate matching capacity is \(4(M_{\vartheta,-}+M_{\vartheta,+})\). By \eqref{eq:legacy-K.9}, \(J_\vartheta\) decreases faster in absolute units than this capacity, so their difference \(J_\vartheta a_{+,\vartheta}\) declines. The matching capacity nevertheless falls faster in proportion than \(J_\vartheta\) throughout \(0<\vartheta<1/2\), leaving a larger relative share for the background coordinate.

Figure~\ref{fig:flow-shares} compares two examples: independence (\(\vartheta=0\), model U) and \(\vartheta=1/4\) (model D). Comparing the triples \((J_\vartheta,a_{+,\vartheta},J_\vartheta a_{+,\vartheta})\), namely \((5/32,13/15,13/96)\) and \((143/960,379/429,379/2880)\), we find that the minimum background share rises from about \(0.8667\) to \(0.8834\), while its unnormalized displacement moment falls from about \(0.1354\) to \(0.1316\).

\begin{figure}[tbp]
\centering
\includegraphics[width=\linewidth]{figures/optimal-flow-shares.pdf}
\caption{An optimal allocation of normal displacement on the tail
surface \(D\), at the equal-cost boundary, for
\(\vartheta=0\) (U, top) and \(\vartheta=1/4\) (D, bottom).
The panels evaluate the explicit shares \(a,b,1-a-b\) in
\eqref{eq:legacy-K.7}; all use the same scale from zero to one.
On each same-side section we allocate displacement to a rebalanced
coordinate and match its projection using the cross-region.
We allocate the remainder to the background coordinate.
Blank regions lie outside \(D\).}
\label{fig:flow-shares}
\end{figure}

\subsection{5.3 Implications for reverse stress testing and further questions}\label{conclusions}\label{v29-concluding-discussion}\label{stress-test-implications}
An exact initial sensitivity need not predict how much dependence stress is required to reverse a diversification trade. The distinction is relevant to reverse sensitivity analysis \cite{pesenti2019,pesenti2022}, where the quantity of interest is the distance to a decision reversal. Under the traded-coordinate dominance condition \eqref{eq:legacy-3.4}, \(r_{0,\varepsilon}=m_\varepsilon/K_\varepsilon\) extrapolates an attained initial slope, yet if \(\gamma_1=c_{\min}<c_3\),
\[
\frac{\rho^\pm_\varepsilon}{r_{0,\varepsilon}}\longrightarrow\frac{s^*}{\gamma_1J}>1.
\]
Proposition~\ref{prop:balanced-linear-profile} attains the initial slope by balancing opposite movements on selected fibers. Equality in the efficiency bound at reversal would require balancing cheapest-coordinate projections across the entire tail section, which the corner obstruction of Theorem~\ref{thm:zero-premium} excludes under the stated cost condition. Unlike optimal allocation across full Fr\'echet or Wasserstein ambiguity sets \cite{liu2026}, our reverse stress test fixes a proposed trade and measures the least perturbation of a common reference law that reverses its paired ES ranking.

For a whole-book comparison across reference dependence structures, the size of the nominal advantage must also enter the assessment. In the family of Section~5.2, consider the boundary perturbation \(\boldsymbol\gamma(\eta)=(\gamma_0,\gamma_0,\gamma_0(1+\eta))\). The gain-normalized critical budget and its right derivative at equal effective costs satisfy
\[
\mathcal R_\vartheta(\eta):=\lim_{\varepsilon\downarrow0}
\frac{\rho^\pm_{\varepsilon,\vartheta}(\eta)}{m_{\varepsilon,\vartheta}}
=\frac{p\,s^*_\vartheta(\eta)}{J_\vartheta},\qquad
\mathcal R_\vartheta'(0+)=p\gamma_0a_{+,\vartheta}.
\]
Theorem~\ref{thm:projected-flow} gives the absolute right derivative \((s^*_\vartheta)'(0+)=\gamma_0J_\vartheta a_{+,\vartheta}\). By \eqref{eq:legacy-7.2}, increasing \(\vartheta\) raises the gain-normalized response \(\mathcal R_\vartheta'(0+)\) but lowers the absolute response: \mbox{\(\frac{d}{d\vartheta}(\gamma_0J_\vartheta a_{+,\vartheta})=\gamma_0(J_\vartheta'a_{+,\vartheta}+J_\vartheta a_{+,\vartheta}')<0\)}. Thus the decline in the nominal-gain scale outweighs the increase in background share. Reporting both boundary sensitivities to a risk committee distinguishes a greater relative need for background displacement from a smaller absolute response in budget units.

Beyond the guaranteed initial linear interval, characterizing the boundaries between the active and stationary regions of an optimal kernel could help determine the full profile \(F\) between the analytical envelopes. A uniform rate for \(\sup_{s\le S_0}|\varepsilon^{-2}D_\varepsilon(s\varepsilon^2)-F(s)|\), combined with the lower growth estimate in Theorem~\ref{thm:strict-crossing}, would yield finite-trade error bounds for the critical budgets.

In the compact-box benchmark, a restoration premium can arise from dependence geometry and full marginal conservation even with bounded marginal tails. For unbounded marginals such as integrable Pareto or lognormal laws, the inverse-CDF derivative is unbounded near probability one, so the uniform displacement estimate in the recovery argument would need a tail-weighted substitute. One-dimensional \(W_1\) remains the \(L^1\) integral of the CDF discrepancy, so this loss of a uniform quantile bound does not by itself imply a different restoration-cost exponent. The open problem is to control that integral jointly with the paired-risk error and determine whether the \(\varepsilon^2\) critical scale persists under suitable tail assumptions. Higher-dimensional projected-flow representations and estimation of critical budgets from loss data raise related questions, alongside recent work on ES estimation under weak tail assumptions \cite{bartleckstein2026}.

\section{Appendix A. Analytic proofs}\label{app:reference-calculus}
\subsection{A.1 Reference ES derivatives}
The first-order formulas are the usual allocation and quantile derivatives \cite{scaillet2004,tasche2001}, and the conditional-variance second derivative appears in \cite[Section~3.2]{martin2011}. We give the section argument to verify that continuity of \(f\) on the compact box suffices without differentiating \(f\).
\resultnumber{A.1}\begin{lemma}[First and second variations on a compact box]\label{lemma:general-variation}\label{lemma:reference-calculus}
Let a reference law have a positive continuous density \(f\) on a product \(K\) of nondegenerate compact intervals in \(\mathbb R^n\). For \(B=b\cdot X\), \(b\ne0\), and \(H=h\cdot X\), let \(v\) be the upper-tail \(p\) threshold of \(B\), and set
\[
m_k(\eta,t)=\int_{K\cap\{B+\eta H=t\}}
\frac{H(x)^k f(x)}{|b+\eta h|}\,d\mathcal H^{n-1}(x),
\qquad k=0,1,2.
\]
For sufficiently small \(\eta\), the threshold \(t(\eta)\) of \(B+\eta H\) is unique and differentiable. Writing \(m_k=m_k(0,v)\) and
\(T_H=E[H1_{\{B>v\}}]\), one has \(m_0>0\) and
\[
\begin{aligned}
t(\eta)&=v+\frac{m_1}{m_0}\eta+o(\eta),\\
ES_p(B+\eta H)-ES_p(B)
&=\frac{T_H}{p}\eta+
\frac1{2p}\left(m_2-\frac{m_1^2}{m_0}\right)\eta^2+o(\eta^2).
\end{aligned}
\]
\end{lemma}
\begin{proof}
Fix \(j\) with \(b_j\ne0\) and eliminate \(x_j\) from \(B+\eta H=t\). The Jacobian \(|b_j+\eta h_j|^{-1}\) agrees with the coarea density because \(d\mathcal H^{n-1}=(|b+\eta h|/|b_j+\eta h_j|)\,dx_{-j}\). To apply dominated convergence as \((\eta,t)\) varies, we need to control where the solved coordinate enters or leaves its interval. If another component of \(b\) is nonzero, it remains nonzero after a sufficiently small perturbation, and at each fixed nearby parameter pair the boundary points lie on at most two affine hyperplanes of zero \(dx_{-j}\)-measure. Along any sequence converging to that pair, the section integrands therefore converge outside this null set. If \(b_j\) is the only nonzero component, the interior level \(v\) instead keeps the solved coordinate uniformly away from both endpoints throughout a small parameter neighborhood. Compactness of the base and the bound away from zero for \(|b_j+\eta h_j|\) provide a common integrable bound in either case, proving continuity of the section moments. Since an interior section contains a relatively open patch and \(f>0\), we also have \(m_0>0\).

For \(Q(\eta,t)=P(B+\eta H\le t)\) and \(\Psi(\eta,t)=t+p^{-1}E(B+\eta H-t)_+\), differentiating the moving endpoint gives \(Q_t=m_0\), \(Q_\eta=-m_1\), \(\Psi_{tt}=m_0/p\), \(\Psi_{\eta t}=-m_1/p\), and \(\Psi_{\eta\eta}=m_2/p\), without differentiating \(f\). The implicit function theorem applied to \(Q(\eta,t(\eta))=1-p\) yields \(t'=m_1/m_0\). Along the optimizing threshold, \(\Psi_t(\eta,t(\eta))=0\), so
\[
\begin{aligned}
\frac{d}{d\eta}ES_p(B+\eta H)&=p^{-1}E[H1_{\{B+\eta H>t(\eta)\}}],\\
\frac{d^2}{d\eta^2}ES_p(B+\eta H)&=\left(\Psi_{\eta\eta}-\frac{\Psi_{\eta t}^2}{\Psi_{tt}}\right)(\eta,t(\eta))
=\frac1p\left(m_2-\frac{m_1^2}{m_0}\right)(\eta,t(\eta)).
\end{aligned}
\]
Taylor expansion at zero gives the stated variation formula.
\end{proof}

\label{appendix-a.-initial-slope-and-tail-surface-capacity}
\subsection{A.2 Threshold localization}\label{proof:threshold-localization}
\resultnumber{A.2}\begin{lemma}[Threshold localization]\label{b-l2-all-optimal-thresholds-under-a-budget-of-order-varepsilon2} For every finite \(M\), there are \(K_M,\varepsilon_M>0\) such that, for \(0<\varepsilon<\varepsilon_M\), any \(P\in\mathcal P(S)\) with \(W_\gamma(P,P_0)\le M\varepsilon^2\) has every ES-minimizing threshold of either portfolio loss in the window
\begin{equation}
|t-v|\le K_M\varepsilon.
\tag{A.1}\label{eq:legacy-B.6}
\end{equation}
\end{lemma}

\textit{Proof.}
For every \(P\) in Lemma~\ref{b-l2-all-optimal-thresholds-under-a-budget-of-order-varepsilon2}, the Lipschitz bound for the portfolio loss gives, uniformly in \(t\),

\begin{equation}
|\mathscr E_{L_{c,\varepsilon}}(P;t)
-\mathscr E_{L_{c,\varepsilon}}(P_0;t)|
\le\frac{\kappa_B+\varepsilon \kappa_H}{p}W_\gamma(P,P_0).
\tag{A.2}\label{eq:legacy-B.7}
\end{equation}

Evaluating \eqref{eq:legacy-B.7} at any minimizer \(t\) under \(P\) and at a reference minimizer bounds the reference suboptimality at \(t\) by twice its right-hand side. By continuity of the section density in Lemma~\ref{lemma:general-variation}, \(m_0(c\varepsilon,t)\ge g/2\) on a fixed neighborhood of \(v\) for all sufficiently small \(\varepsilon\). The unique reference minimizers \(v_{c,\varepsilon}\), with \(v_{0,\varepsilon}=v\) and \(v_{1,\varepsilon}=t_\varepsilon\), converge to \(v\). Integrating the curvature bound within this neighborhood, we obtain
\begin{equation}
\mathscr E_{L_{c,\varepsilon}}(P_0;t)
-ES_p^{P_0}(L_{c,\varepsilon})
\ge\frac{g}{4p}|t-v_{c,\varepsilon}|^2.
\tag{A.3}\label{eq:legacy-B.8}
\end{equation}

Every candidate minimizer lies in the compact loss range. Since \(v_{c,\varepsilon}\to v\), \eqref{eq:legacy-B.8} gives a positive reference suboptimality gap at the endpoints of a smaller fixed neighborhood of \(v\), uniformly for small \(\varepsilon\). Convexity extends this gap outside the neighborhood, contradicting the uniform \(O(\varepsilon^2)\) bound from \eqref{eq:legacy-B.7}. Applying \eqref{eq:legacy-B.8} inside gives \(|t-v_{c,\varepsilon}|=O_M(\varepsilon)\), and \eqref{eq:legacy-3.2} yields \eqref{eq:legacy-B.6} for every minimizing threshold.\unskip\nobreak\hfill\proofbox

\subsection{A.3 Projections and threshold convexity}
\label{a.4-source-fluxes-and-moving-optimal-thresholds}
\label{a.5-strict-source-imbalance-on-an-exchangeable-tail-surface}
For a finite-\(I\) kernel, the conditional mean \(e_i(\xi,u)=\int d_i\alpha_{\xi,u}(dd)\) belongs to \(L^1(\nu)\), so Fubini's theorem gives

\begin{equation}
j_1(q)=\int_{\mathbb R}\int_I
\frac{f(q,y,z_v(q,y))}{\beta}1_\Omega(q,y)
e_1((q,y,z_v(q,y)),u)\,dy\,du.
\tag{A.4}\label{eq:legacy-B.13}
\end{equation}

Interchanging \(x,y\) gives the second-coordinate formula. For the third projection we parameterize the slice by \(y=v-\beta q-x\). The area Jacobian \(\beta\) cancels the factor \(1/\beta\) in \(\sigma\), yielding

\begin{equation}
j_3(q)=\int_{\mathbb R}\int_I
f(x,v-\beta q-x,q)1_{\{\ell<v-\beta q-x<r\}}
e_3((x,v-\beta q-x,q),u)\,dx\,du.
\tag{A.5}\label{eq:legacy-B.14}
\end{equation}

Together with the second-coordinate formula, these identities give \(j_i\in L^1(I_i)\) and

\begin{equation}
\sum_i\gamma_i\|j_i\|_1\le I,\qquad I\le\mathcal C\le2I.
\tag{A.6}\label{eq:legacy-B.15}
\end{equation}

The local plan need not balance internally: \(\int_{I_i}j_i(q)\,dq=\int d_i\,d\Pi_\alpha\) may be nonzero, with compensation supplied by the original plan outside the retained layer. For the directional moments \(\ell_-,\ell_+\) in \eqref{eq:legacy-A.19}, exchangeability equates their totals but not their distributions along the first coordinate.

\resultnumber{A.3}\begin{lemma}[Positive imbalance on an exchangeable tail surface]\label{lem:projected-imbalance} The quantities in \eqref{eq:legacy-A.19},

\begin{equation}
J=\int_I(\ell_-+\ell_+)\,dx,\qquad
\chi=\frac{\int_I|\ell_--\ell_+|\,dx}{J}
\tag{A.7}\label{eq:legacy-A.24}
\end{equation}

satisfy \(J>0\) and \(0<\chi<1\), where \(J\) is the section moment in \eqref{eq:legacy-3.1}.
\end{lemma}

\begin{proof}
Because exchangeability balances the total downward and upward moments, we test their difference against the coordinate \(x\) to detect fiberwise imbalance. Set
\begin{equation*}
M_j=\int_I\int_I|x-y|^j k(x,y)\,dy\,dx,\qquad
\mathfrak b(x)=\ell_-(x)-\ell_+(x).
\end{equation*}
Then \(J=M_2/(2\beta)\) and
\begin{equation*}
\mathfrak b(x)=\frac1{2\beta}\int(x-y)|x-y|k(x,y)\,dy.
\end{equation*}
Antisymmetry gives zero total mass, but testing against the coordinate and using \(k(x,y)=k(y,x)\) yields
\begin{equation}
\int \mathfrak b(x)\,dx=0,\qquad
\int x\,\mathfrak b(x)\,dx=\frac{M_3}{4\beta}.
\tag{A.8}\label{eq:legacy-A.25}
\end{equation}
Centering the coordinate at \(m_I=(\ell+r)/2\) gives \(|x-m_I|\le \Delta_I/2\), and therefore
\begin{equation}
\int_I|\mathfrak b|\ge\frac2{\Delta_I}\int_I(x-m_I)\mathfrak b(x)\,dx=\frac{M_3}{2\beta \Delta_I},
\qquad
\chi\ge\frac{M_3}{\Delta_IM_2}>0.
\tag{A.9}\label{eq:legacy-A.26}
\end{equation}
The interior diagonal point from \hyperref[geometry:interior-diagonal]{Section~2.3} has a neighborhood in \(\Omega\) where \(f>0\). Its off-diagonal part has positive area, so \(M_2,M_3>0\), proving \(J>0\) and \(\chi>0\).

On a smaller central interval of \(x\)-values, each side \(y<x\) and \(y>x\) still contains a positive-length section inside \(\Omega\). Both \(\ell_-(x)\) and \(\ell_+(x)\) are therefore positive, giving
\begin{equation}
2\int\min(\ell_-,\ell_+)=J-\int|\ell_--\ell_+|=(1-\chi)J>0,
\tag{A.10}\label{eq:legacy-A.27}
\end{equation}
and hence \(\chi<1\).
\end{proof}

\keepproofstart{9\baselineskip}
\paragraph{Convexity of limiting threshold objectives}
\label{b.0-model-objects-and-main-statement}
For finite-cost kernels, convexity of \(G_c\) follows by identifying its distributional second derivative with the measure of transported rescaled losses.

\resultnumber{A.4}\begin{lemma}[Threshold objectives of finite-cost kernels]\label{lem:threshold-convexity} For every kernel with finite \(I\) and every \(c\in\{0,1\}\), the transported rescaled loss \(u+b\cdot d+cH\) has a locally finite pushforward measure \(\lambda_c\), and

\begin{equation}
(G_c^\alpha)''=\lambda_c\ge0
\quad\text{in the sense of distributions.}
\tag{A.11}\label{eq:legacy-B.Z.1}
\end{equation}

Thus both threshold objectives are continuous and convex and attain their minima. For each objective, zero lying between its one-sided derivatives characterizes a globally optimal threshold.
\end{lemma}

\begin{proof}
Fix \([-T,T]\) and \(R_0>T+\Delta_I\). The source mass on \(|u|\le R_0\) is finite, and any source outside it reaching this loss interval must satisfy \(\kappa_B|d|_\gamma\ge R_0-T-\Delta_I\). Its mass is at most \(\kappa_BI/(R_0-T-\Delta_I)\), proving local finiteness of \(\lambda_c\). Translation invariance in \(u\) gives the source rescaled-loss measure \(g\,dt\).

For a smooth compactly supported test function \(\varphi\), we use \(\int(z-\tau)_+\varphi''(\tau)d\tau=\varphi(z)\). The bound \(|\mathfrak q_c|\le \kappa_B|d|_\gamma\) provides the integrable majorant \(\kappa_B|d|_\gamma\|\varphi''\|_1\), so Fubini's theorem gives
\begin{equation}
\langle (T_c^\alpha)'',\varphi\rangle
=\int[\varphi(u+b\cdot d+cH)-\varphi(u+cH)]\,d\Pi_\alpha.
\tag{A.12}\label{eq:legacy-B.Z.2}
\end{equation}
We obtain \((T_c^\alpha)''=\lambda_c-g\,dt\), so adding the quadratic derivative proves \eqref{eq:legacy-B.Z.1}. To conclude, dominated convergence of the hinge increments gives continuity, and \(|T_c|\le\kappa_BI\) gives coercivity of \(G_c\). Continuity and coercivity ensure attainment, and convexity gives the stated one-sided derivative criterion.
\end{proof}

\keepproofstart{9\baselineskip}
\subsection{A.4 Preliminary laws and kernel approximation}\label{proof:preliminary-cdf}
The recovery class \(\mathcal K_{\rm fin}^c\) introduced after Theorem~\ref{thm:critical-limit} has the explicit form, for finitely many fixed vectors \(d_k\),
\begin{equation}
\begin{aligned}
&\alpha_k\in C_c(\Omega\times\mathbb R),\quad
\alpha_k\ge0\ (1\le k\le N),\quad\sum_{k=1}^N\alpha_k\le1,\\
&\alpha_{\xi,u}=\sum_{k=1}^N\alpha_k(\xi,u)\delta_{d_k}
+\left(1-\sum_{k=1}^N\alpha_k(\xi,u)\right)\delta_0.
\end{aligned}
\tag{A.13}\label{eq:legacy-B.26}
\end{equation}
For \(Z_\varepsilon=(x,y,z_v(x,y)+\varepsilon u/\beta)\) and \(f_\varepsilon=f(Z_\varepsilon)\), translate branch \(k\) to \(Z_\varepsilon+\varepsilon d_k\) and leave the remaining mass fixed. Compact interior activity keeps these sources and targets in \(S\) for small \(\varepsilon\), and translation of each subdensity gives an absolutely continuous law \(Q_\varepsilon\). The source volume is \(\varepsilon\,dx\,dy\,du/\beta\).

Let \(F_i\) be the reference marginal CDF and \(F_{i,\varepsilon}\) that of \(Q_\varepsilon\). The fluxes are continuous with compact support, and we claim
\begin{equation}
\sup_{q\in I_i}
\left|\frac{F_{i,\varepsilon}(q)-F_i(q)}{\varepsilon^2}+j_i(q)\right|\to0.
\tag{A.14}\label{eq:legacy-B.28}
\end{equation}
For each branch, let \(a_{ik,\varepsilon}\) be the density of
\((Z_{\varepsilon,i})_\#(\alpha_k f_\varepsilon\,dx\,dy\,du/\beta)\).
The CDF change is the integral over the oriented interval crossed by that branch, hence
\begin{equation}
\frac{F_{i,\varepsilon}(q)-F_i(q)}{\varepsilon^2}
=-\sum_kd_{ki}\int_0^1a_{ik,\varepsilon}(q-\varepsilon t d_{ki})\,dt.
\tag{A.15}\label{eq:legacy-B.29}
\end{equation}
For \(i=1,2\), the densities are obtained by integrating \(\alpha_k f_\varepsilon/\beta\) over the other surface coordinate and \(u\). For \(i=3\), we use coordinates \((w,y)\), with \(w=z_v\), whose area Jacobian is \(\beta\). Since \(Z_{\varepsilon,3}=w+\varepsilon u/\beta\), evaluating the third-coordinate marginal at \(q-\varepsilon t d_{k3}\) in \eqref{eq:legacy-B.29} requires
\begin{equation}
w=q-\varepsilon u/\beta-\varepsilon t d_{k3},
\tag{A.16}\label{eq:legacy-B.30}
\end{equation}
with weight \(\alpha_k(v-y-\beta w,y,u)f(v-y-\beta w,y,w+\varepsilon u/\beta)\), after cancellation of the two Jacobian factors. Compact interior support and uniform continuity give \(a_{ik,\varepsilon}\to a_{ik,0}\) uniformly, including the zero extensions at the endpoints. Since \(j_i=\sum_kd_{ki}a_{ik,0}\), \eqref{eq:legacy-B.29} proves \eqref{eq:legacy-B.28} for all three coordinates.

\paragraph{Approximation by finite-branch kernels}
\label{b-l8-bounded-measurable-kernels-arbitrary-finite-i-kernels-and-diagonal-recovery}
To approximate a general finite-\(I\) kernel by kernels in \(\mathcal K_{\rm fin}^c\), we first retain activity on \(A_n=\{\xi\in\Omega_n,\ |u|\le n,\ |d|_\gamma\le n\}\), where \(\Omega_n\uparrow\Omega\) is a compact exhaustion, and return the omitted conditional probability to \(d=0\). With \(I_{\mathrm{tail},n}=\int_{A_n^c}|d|_\gamma\,d\Pi_\alpha\to0\), projection contraction and the displacement-Lipschitz bound for hinge increments give
\begin{equation}
|I_n-I|\le I_{\mathrm{tail},n},\qquad
\sum_i\gamma_i\|j_i^n-j_i\|_1\le I_{\mathrm{tail},n},\qquad
\sup_\tau|T_c^{\alpha_n}-T_c^\alpha|\le \kappa_BI_{\mathrm{tail},n}.
\tag{A.17}\label{eq:legacy-B.40}
\end{equation}
In particular, \(\mathcal C(\alpha_n)\to\mathcal C(\alpha)\) and \(\mathcal V(\alpha_n)\to\mathcal V(\alpha)\). The loss of first moment bounds any increase of the projected norms, so also \(\mathcal C(\alpha_n)\le\mathcal C(\alpha)\).

On a fixed truncated window with \(|u|\le U\), replace displacement cells of diameter at most \(\eta\) by representatives \(d_k\). Their conditional masses are measurable weights \(a_k\ge0\), \(\sum_k a_k\le1\), and the same bounds yield
\begin{equation}
|I_\eta-I|+\sum_i\gamma_i\|j_{i,\eta}-j_i\|_1
+\sum_{c=0}^1\sup_\tau|T_{c,\eta}(\tau)-T_c(\tau)|
\le C\eta\,\nu(\Omega\times[-U,U]).
\tag{A.18}\label{eq:legacy-B.38}
\end{equation}
To regularize the weights while preserving their simplex constraint, we extend them by zero, convolve with a common nonnegative approximate identity, and multiply by a common compact interior cutoff taking values in \([0,1]\). Since the section density is bounded, we have \(\sum_k\|a_k^{(m)}-a_k\|_{L^1(\nu)}\to0\). With the finite set of displacements fixed in this step, this convergence controls \(I\), all projected fluxes in \(L^1\), and \(T_0,T_1\) uniformly in the threshold. For a compactly supported Lipschitz test on \(Z\), the value at \(d=0\) contributes the same source-only integral to the original and approximating plans and cancels in their difference, leaving errors controlled by the preceding truncation, quantization, and smoothing bounds. Density of these tests on compact sets and successive choice of the three approximation parameters yield kernels \eqref{eq:legacy-B.26} converging in \(I,\mathcal C,\mathcal V\) and the vague plan topology.


\subsection{A.5 Proof of the variational limit}\label{v21-limit-proof}\label{appendix-b-proof-of-the-critical-variational-limit-under-full-marginal-constraints}

\begin{proof}[Proof of Theorem~\ref{thm:critical-limit}]
\phantomsection\label{b-l3-local-compactness-of-arbitrary-transport-plans-and-the-three-fluxes}
We first establish vague compactness and the energy lower bound in part~(i). Let \(\pi_n\) be admissible plans from \(P_0\) to \(P_n\in\Pi\) along a subsequence realizing a finite energy liminf, with rescaled coordinates and lifts
\begin{equation}
u=\frac{B(X)-v}{\varepsilon_n},\qquad
d=\frac{Y-X}{\varepsilon_n},\qquad
\Pi_n=\mathscr T_{\varepsilon_n}\pi_n,
\tag{A.19}\label{eq:legacy-B.9}
\end{equation}
and the bounded energies give
\begin{equation}
\int|d|_\gamma d\Pi_n=\mathcal E_{\varepsilon_n}(\Pi_n)\le M.
\tag{A.20}\label{eq:legacy-B.11}
\end{equation}
\phantomsection\label{proof:full-source-kernel}
Since \(X=(x,y,z_v(x,y)+\varepsilon_nu/\beta)\), their source densities are
\begin{equation}
w_n(x,y,u)=\frac1\beta f(x,y,z_v(x,y)+\varepsilon_nu/\beta)
1_{\{\ell_3\le z_v(x,y)+\varepsilon_nu/\beta\le r_3\}}.
\tag{A.21}\label{eq:legacy-B.10}
\end{equation}
By dominated convergence and the zero area of the section boundary, \(w_n\to f(\xi)1_\Omega/\beta\) in \(L^1\) on bounded \(u\)-windows. Because the source masses are uniformly bounded on each such window and \eqref{eq:legacy-B.11} bounds mass with \(|d|_\gamma>K\) by \(M/K\), we can extract a vague subsequential limit through compact exhaustion. To identify its full source as \(\nu=\sigma\,du\), we first test the source projection with a displacement cutoff and then let \(K\to\infty\). The limit consequently disintegrates into probability kernels \(\Pi_\alpha=\nu\alpha\), with
\begin{equation}
I(\alpha)\le\liminf_n\mathcal E_{\varepsilon_n}(\Pi_n).
\tag{A.22}\label{eq:legacy-B.12}
\end{equation}
By Appendix A.3, the projected moments have densities \(j_i\in L^1(I_i)\). To include their marginal-restoration cost in the lower bound, we choose continuous cutoffs \(0\le\chi_R\uparrow1\), compactly supported in the interior surface and bounded \(u,d\)-windows.
\phantomsection\label{b-l4-the-compensation-cost-forced-by-full-marginal-constraints}
The coordinate CDF discrepancy of \(\pi_n^R=\chi_R\pi_n\) is
\begin{equation}
\Delta_{i,n}^R(q)=\int[1_{\{Y_i\le q\}}-1_{\{X_i\le q\}}]\,d\pi_n^R.
\tag{A.23}\label{eq:legacy-B.16}
\end{equation}
Because the full plan preserves marginal \(i\), the complementary nonnegative plan has discrepancy \(-\Delta_{i,n}^R\). Its cost is bounded below by the one-dimensional crossing identity: for any nonnegative plan \(\lambda\),
\[
\int|Y_i-X_i|\,d\lambda
=\int_{I_i}\int|1_{\{Y_i\le q\}}-1_{\{X_i\le q\}}|\,d\lambda\,dq
\ge\int_{I_i}\left|\int(1_{\{Y_i\le q\}}-1_{\{X_i\le q\}})\,d\lambda\right|dq.
\]
Applying this to the complement and adding the local cost gives
\begin{equation}
\varepsilon_n^2\mathcal E_{\varepsilon_n}(\Pi_n)\ge
\int|Y-X|_\gamma\,d\pi_n^R
+\sum_i\gamma_i\int_{I_i}|\Delta_{i,n}^R(q)|\,dq.
\tag{A.24}\label{eq:legacy-B.17}
\end{equation}
For \(\varphi\in C(I_i)\), integration over the oriented interval from \(X_i\) to \(Y_i\) gives
\begin{equation}
\varepsilon_n^{-2}\int\varphi(q)\Delta_{i,n}^R(q)\,dq
=-\int\chi_Rd_i\int_0^1\varphi(X_i+t\varepsilon_nd_i)\,dt\,d\Pi_n.
\tag{A.25}\label{eq:legacy-B.18}
\end{equation}
On the cutoff support, \(X\to\xi\) uniformly and \(d\) is bounded. Thus local weak convergence identifies the projected moment \(j_i^Rdq=(\xi_i)_\#(\chi_Rd_i\,d\Pi_\alpha)\) through
\begin{equation}
\varepsilon_n^{-2}\Delta_{i,n}^R(q)\,dq
\rightharpoonup-j_i^R(q)\,dq
\quad\hbox{as finite signed measures.}
\tag{A.26}\label{eq:legacy-B.19}
\end{equation}
Lower semicontinuity of total variation in \eqref{eq:legacy-B.17} now yields \(\liminf_n\mathcal E_{\varepsilon_n}(\Pi_n)\ge I_R+\sum_i\gamma_i\|j_i^R\|_1\), where \(I_R=\int\chi_R|d|_\gamma d\Pi_\alpha\). We remove the cutoff using
\begin{equation}
I_R\uparrow I,\qquad
\sum_i\gamma_i\|j_i-j_i^R\|_1
\le\int(1-\chi_R)|d|_\gamma d\Pi_\alpha\to0.
\tag{A.27}\label{eq:legacy-B.20}
\end{equation}
Thus any vaguely convergent sequence \(\Lambda_n\to\Lambda\) with finite energy liminf has \(\Lambda=\Pi_\alpha\) and \(\mathcal C(\alpha)\le\liminf_n\mathcal E_{\varepsilon_n}(\Lambda_n)\). In particular, the liminf must be infinite when \(\mathcal E(\Lambda)=+\infty\), proving the \(\Gamma\)-lower bound on all of \(\mathscr X\). Choosing optimal plans gives the Wasserstein bound in \eqref{eq:legacy-4.5}.

\phantomsection\label{lem:v21-coordinate-lsc}\label{l1.-lower-semicontinuity-of-the-combined-cost-of-one-coordinate}
For later boundary sensitivity, we also need lower semicontinuity of \(A_i(\alpha)=M_i(\alpha)+\|j_i\|_1\), where \(M_i(\alpha)=\int|d_i|\,d\Pi_\alpha\). Retaining activity weighted by \(\chi_R\) and assigning the remainder to \(d=0\) defines \(\alpha^R\), for which projection contraction gives
\begin{equation}
\|j_i-j_i^R\|_1\le M_i-M_i^R,\qquad
A_i(\alpha^R)\le A_i(\alpha),\qquad A_i(\alpha^R)\to A_i(\alpha).
\tag{A.28}\label{eq:legacy-H.5}
\end{equation}
For fixed \(R\), local weak convergence implies convergence of \(M_i(\alpha_n^R)\) and every projected test \(\int\varphi(\xi_i)d_i\chi_R\,d\Pi_{\alpha_n}\). Lower semicontinuity of total variation followed by \eqref{eq:legacy-H.5} as \(R\to\infty\) gives
\begin{equation}
A_i(\alpha)\le\liminf_n A_i(\alpha_n),\qquad
\mathcal C(\alpha)\le\liminf_n\mathcal C(\alpha_n).
\tag{A.29}\label{eq:legacy-H.7}
\end{equation}
The same extraction with fixed source \(\nu\) makes cost sublevels compact in the vague topology, that is, convergence against \(C_c(Z)\). A countable compact exhaustion metrizes this topology on the sublevels.

\paragraph{The risk limit}\label{b-l5-renormalization-of-the-common-term-and-identification-of-the-risk-difference}

Although \(|\mathfrak q_c|\le\kappa_B|d|_\gamma\) makes each hinge integral finite, bounded moments alone do not imply uniform integrability. We therefore subtract the common increment and estimate \(r_c(\tau)=\mathfrak q_c(\tau)-\mathfrak q_0(0)\). For \(w=b\cdot d\), the function \(z\mapsto(z+w)_+-z_+\) is 1-Lipschitz, bounded in magnitude by \(|w|\), and constant on each exterior half-line determined by its breakpoints \(0\) and \(-w\). The two arguments defining \(r_c\), namely \(u+cH-\tau\) and \(u\), differ by at most \(A_T=\|H\|_\infty+T\) when \(|\tau|\le T\). If \(|u|>A_T+|w|\), both arguments lie on the same exterior half-line, so their increments agree. These observations give
\begin{equation}
|r_c(\tau)|\le A_T,\quad |r_c(\tau)|\le2|b\cdot d|,
\quad r_c(\tau)\ne0\Longrightarrow |u|\le A_T+|b\cdot d|.
\tag{A.30}\label{eq:legacy-B.21}
\end{equation}
When \(R>A_T\), a nonzero remainder outside \(\{|u|\le R,|d|_\gamma\le R\}\) requires \(|d|_\gamma>k_R=\min\{R,(R-A_T)/\kappa_B\}\). The first-moment bound for both the sequence and its limit gives
\begin{equation}
\int_{\{|u|>R\}\cup\{|d|_\gamma>R\}}|r_c(\tau)|\,d\Pi_n
\le A_TM/k_R\longrightarrow0.
\tag{A.31}\label{eq:legacy-B.22}
\end{equation}
On bounded windows, vague convergence and the common threshold-Lipschitz bound give uniform convergence in \(\tau\). Because the surface-boundary cutoff has vanishing source mass, we remove it first and then remove the other cutoffs using \eqref{eq:legacy-B.22}, obtaining
\begin{equation}
\int r_c(\tau)d\Pi_n\longrightarrow
T_c^\alpha(\tau)-T_0^\alpha(0)
\quad\hbox{uniformly on every bounded }\tau\hbox{ interval.}
\tag{A.32}\label{eq:legacy-B.23}
\end{equation}

The exact target loss also contains \(c\varepsilon_n^2h\cdot d\), with \(h=(1,-1,0)\), whose rescaled integrated contribution is at most \(\varepsilon_n\kappa_H\int|d|_\gamma d\Pi_n=O(\varepsilon_n)\). Combining the transport identity with \(H(X)=H(\xi)=x-y\) and the uniform reference expansion \eqref{eq:legacy-B.2}, and writing \(Q_n=\int\mathfrak q_0(0)d\Pi_n\), gives
\begin{equation}
\frac p{\varepsilon_n^2}
\{\mathscr E_{L_{c,\varepsilon_n}}(P_n;v+\varepsilon_n\tau)
-ES_p^{P_0}(B)\}-Q_n
=\frac g2\tau^2+c^2J+\int r_c(\tau)d\Pi_n+o(1).
\tag{A.33}\label{eq:legacy-B.24}
\end{equation}
Lemma~\ref{b-l2-all-optimal-thresholds-under-a-budget-of-order-varepsilon2} and \(g\tau^2/2\le2\kappa_BI\) place all original and limiting minimizers in a common compact interval. The finite scalar \(Q_n\), with \(|Q_n|\le\kappa_BM\), translates both objectives equally and cancels from their difference of minima without needing to converge itself. Uniform convergence on that interval therefore yields
\begin{equation}
\Phi_{\varepsilon_n}(P_n)/\varepsilon_n^2\to\mathcal V(\alpha).
\tag{A.34}\label{eq:legacy-B.25}
\end{equation}
Since plan cost bounds target distance, localization applies to every bounded-energy sequence in part (i). Optimal plans give the corresponding law-level limits in \eqref{eq:legacy-4.5}.

\paragraph{Exact recovery by quantile restoration}\label{b-l6-nonnegative-preliminary-recovery-for-finite-smooth-kernels-and-a-uniform-cdf-expansion}

For a fixed kernel \eqref{eq:legacy-B.26}, the preliminary plan of Appendix A.4 modifies \(O_\alpha(\varepsilon)\) source mass and has cost
\begin{equation}
\operatorname{cost}(P_0\to Q_\varepsilon)
=\varepsilon^2\sum_k\int|d_k|_\gamma\alpha_k
\frac{f_\varepsilon}{\beta}\,dx\,dy\,du
=\varepsilon^2 I(\alpha)+o(\varepsilon^2).
\tag{A.35}\label{eq:legacy-B.27}
\end{equation}
Threshold localization applies to \(Q_\varepsilon\) by this cost bound. At a translated source,
\(L_{c,\varepsilon}(Z_\varepsilon+\varepsilon d_k)=v+\varepsilon(u+b\cdot d_k+cH)+c\varepsilon^2h\cdot d_k\).
The layer volume \(\varepsilon\,dx\,dy\,du/\beta\), hinge homogeneity, and the reference expansion give
\begin{equation}
\frac p{\varepsilon^2}
\{\mathscr E_{L_{c,\varepsilon}}(Q_\varepsilon;v+\varepsilon\tau)
-ES_p^{P_0}(B)\}
=c^2J+\frac g2\tau^2+T_c^\alpha(\tau)+o(1).
\tag{A.36}\label{eq:legacy-B.31}
\end{equation}
The remainder is uniform on bounded threshold intervals, by continuity of \(f\) and boundedness of the branch displacements. Minimization on the common compact window gives
\begin{equation}
\Phi_\varepsilon(Q_\varepsilon)/\varepsilon^2\to\mathcal V(\alpha).
\tag{A.37}\label{eq:legacy-B.32}
\end{equation}

\phantomsection\label{b-l7-exact-repair-of-all-three-marginals-repair-cost-and-cancellation-in-the-risk-difference}
To restore the marginals, observe that their reference densities are bounded below by \(m_f=(\min_Sf)\min\{\Delta_I\Delta_3,\Delta_I^2\}>0\). Thus \(F_i^{-1}\) is \(m_f^{-1}\)-Lipschitz. Set
\begin{equation}
R_{i,\varepsilon}(q)=F_i^{-1}(F_{i,\varepsilon}(q)),\quad
R_\varepsilon(z)=(R_{1,\varepsilon}(z_1),R_{2,\varepsilon}(z_2),R_{3,\varepsilon}(z_3)),
\quad P_\varepsilon=(R_\varepsilon)_\#Q_\varepsilon.
\tag{A.38}\label{eq:legacy-B.33}
\end{equation}
The preliminary CDFs are continuous, so \(F_{i,\varepsilon}(Z_i)\) is uniform for \(Z\sim Q_\varepsilon\) and every restored marginal is exactly \(\mu_i\). The uniform CDF expansion \eqref{eq:legacy-B.28} also gives
\begin{equation}
\max_i\sup_{q\in I_i}|R_{i,\varepsilon}(q)-q|=O_\alpha(\varepsilon^2).
\tag{A.39}\label{eq:legacy-B.34}
\end{equation}
Because the monotone coordinate coupling is optimal in one dimension \cite[Chapter~2]{santambrogio2015}, Vallender's formula \cite{vallender1974} identifies its exact restoration cost:
\begin{equation}
\begin{aligned}
E_{Q_\varepsilon}|R_\varepsilon Z-Z|_\gamma
&=\sum_i\gamma_i\int_{I_i}|F_{i,\varepsilon}(q)-F_i(q)|\,dq\\
&=\varepsilon^2\sum_i\gamma_i\|j_i\|_1+o(\varepsilon^2).
\end{aligned}
\tag{A.40}\label{eq:legacy-B.35}
\end{equation}
To show that restoration has no leading paired-risk effect, we estimate its effect uniformly over both minimizing thresholds. The preliminary and restored laws \(Q_\varepsilon,P_\varepsilon\) are within \(O_\alpha(\varepsilon^2)\) of \(P_0\), so Lemma~\ref{b-l2-all-optimal-thresholds-under-a-budget-of-order-varepsilon2} places all their minimizing thresholds in \([v-T_\alpha\varepsilon,v+T_\alpha\varepsilon]\). For any pair in this window, put \(\psi_{t_B,t_A}(z)=(B(z)-t_B)_+-(B(z)+\varepsilon H(z)-t_A)_+\). Away from a sufficiently wide band \(\mathcal B_\varepsilon=\{|B-v|\le K\varepsilon\}\), both hinges stay on the same affine branch along the restoration segment: the gradient of their difference is \(-\varepsilon h\) above both thresholds and zero below both. Inside the band we retain the uniform gradient bound, but only \(O_\alpha(\varepsilon)\) mass is affected. Indeed, the reference density of \(B\) is bounded by \(\Delta_I^2\max_Sf/\beta\) and the preliminary plan changes \(O_\alpha(\varepsilon)\) mass, so \(Q_\varepsilon(\mathcal B_\varepsilon)=O_\alpha(\varepsilon)\). Combining this small band mass and the extra factor \(\varepsilon\) outside the band with the \(O_\alpha(\varepsilon^2)\) restoration displacement in \eqref{eq:legacy-B.34} gives
\begin{equation}
\begin{aligned}
|(E_{P_\varepsilon}-E_{Q_\varepsilon})\psi_{t_B,t_A}|
&\le\varepsilon\kappa_HE_{Q_\varepsilon}|R_\varepsilon Z-Z|_\gamma\\
&\quad+O_\alpha(\varepsilon^2)Q_\varepsilon(\mathcal B_\varepsilon)
=O_\alpha(\varepsilon^3).
\end{aligned}
\tag{A.41}\label{eq:legacy-B.36}
\end{equation}
Choose \(t_B\) optimal under \(Q_\varepsilon\) and \(t_A\) under \(P_\varepsilon\). Separate minimization gives
\[
\Phi_\varepsilon(P_\varepsilon)-\Phi_\varepsilon(Q_\varepsilon)
\le p^{-1}(E_{P_\varepsilon}-E_{Q_\varepsilon})\psi_{t_B,t_A}.
\]
The reversed choice gives the opposite inequality, hence
\begin{equation}
|\Phi_\varepsilon(P_\varepsilon)-\Phi_\varepsilon(Q_\varepsilon)|
=O_\alpha(\varepsilon^3).
\tag{A.42}\label{eq:legacy-B.37}
\end{equation}

With \(\widetilde\pi_\varepsilon\) denoting the preliminary translation plan, the coupling
\[
\bar\pi_\varepsilon=(\operatorname{id}_S,R_\varepsilon)_\#\widetilde\pi_\varepsilon
\in\Gamma(P_0,P_\varepsilon)
\]
has cost at most the sum of \eqref{eq:legacy-B.27} and \eqref{eq:legacy-B.35}. Restoration leaves its source coordinates fixed and changes its lifted displacement by \(O_\alpha(\varepsilon)\), so the bounded source mass on compact windows and uniform continuity of test functions preserve the vague limit \(\Pi_\alpha\). Equations \eqref{eq:legacy-B.32} and \eqref{eq:legacy-B.37} give the paired-risk limit.

For a general finite-\(I\) kernel, choose the approximants \(\alpha^{(m)}\) of Appendix A.4 and a metric \(d_{\rm v}\) for vague convergence. The fixed-kernel construction permits \(\varepsilon_m\downarrow0\) such that, for every \(0<\varepsilon\le\varepsilon_m\),
\[
\begin{aligned}
\mathcal E_\varepsilon(\mathscr T_\varepsilon\bar\pi_\varepsilon^{(m)})
&\le\mathcal C(\alpha^{(m)})+m^{-1},\\
\left|\varepsilon^{-2}\Phi_\varepsilon(P_\varepsilon^{(m)})-\mathcal V(\alpha^{(m)})\right|
+d_{\rm v}(\mathscr T_\varepsilon\bar\pi_\varepsilon^{(m)},\Pi_{\alpha^{(m)}})
&\le m^{-1}.
\end{aligned}
\]
Using the \(m\)-th recovery for \(\varepsilon\in(\varepsilon_{m+1},\varepsilon_m]\) gives a full recovery family with exact marginals, vague limit \(\Pi_\alpha\), and the energy and risk limits in \eqref{eq:legacy-4.6}. Its risk error is \(o(\varepsilon^2)\), because the fixed-approximant cubic constants need not be uniform in \(m\).

\paragraph{Optimizing over the budget}\label{completion-of-b-t-original-model-budgets-budget-boundaries-and-local-uniform-convergence}

\label{b-l9-compact-cost-sublevel-sets-continuity-of-the-risk-difference-and-budget-attainment}
Parts (i)--(ii) are now proved. For kernels with fixed source \(\nu\) and \(I\le M\), the same estimates \eqref{eq:legacy-B.21}--\eqref{eq:legacy-B.23} give uniform convergence of the relative objectives on a common compact interval containing every minimizer. Their common term cancels, so \(\mathcal V\) is continuous on compact cost sublevels and \(F\) is finite and attained, with \(F(0)=-J/p\). Mixing with the identity, \(\alpha^\theta=\theta\alpha+(1-\theta)\delta_0\), scales \(I,j_i,T_c\) linearly and gives
\begin{equation}
\mathcal C(\alpha^\theta)=\theta\mathcal C(\alpha),\quad
|\mathcal V(\alpha^\theta)-\mathcal V(\alpha)|
\le\frac{2\kappa_BI(\alpha)}p(1-\theta).
\tag{A.43}\label{eq:legacy-B.44}
\end{equation}
Applied to a maximizer at \(t>s\), with \(\theta=s/t\), this yields
\begin{equation}
0\le F(t)-F(s)\le\frac{2\kappa_B}{p}(t-s),\qquad0\le s<t.
\tag{A.44}\label{eq:legacy-B.45}
\end{equation}

\paragraph{Reflection at finite cost}\label{sec:reflection-construction}\label{proof:positive-witness}
To obtain a positive value of \(F\) at finite cost, we reflect two normal bands over exchanged surface patches. The \(B\)-increments then cancel, and reversal of \(H\) makes the \(A_\varepsilon\)-increments add with negative sign without shifting the minimizing thresholds. Choose a compact positive-area patch \(E_+\Subset\Omega\cap\{H>h_0\}\), \(h_0>0\), and its exchanged image \(E_-\), of common mass \(m>0\). Set \(\Delta_E=\int_{E_+}H\,d\sigma-\int_{E_-}H\,d\sigma>0\). With \(U>\Delta_I\) and \(R>0\), use \(d=-2ue_1\) on \(E_+\times[U,U+R]\) and \(E_-\times[-U-R,-U]\), leaving other sources fixed. Integration gives
\begin{equation}
I=4\gamma_1m(UR+R^2/2),\qquad \mathcal C\le2I<\infty.
\tag{A.45}\label{eq:legacy-B.46}
\end{equation}
Pairing the positive and negative bands cancels their normal-coordinate and threshold terms, giving, for \(|\tau|<U-\Delta_I\),
\begin{equation}
T_0(\tau)=0,\qquad T_1(\tau)=-\Delta_E R.
\tag{A.46}\label{eq:legacy-B.47}
\end{equation}
The preceding formulas make both objectives stationary at zero, and Lemma~\ref{lem:threshold-convexity} makes them convex. Zero therefore minimizes both objectives, so choosing \(R>J/\Delta_E\) gives the positive paired value
\begin{equation}
\mathcal V=(-J+\Delta_E R)/p>0.
\tag{A.47}\label{eq:legacy-B.48}
\end{equation}

For \(s>0\), let \(\alpha\) maximize \(F(s)\). The mixture \(\alpha^\theta\), \(0\le\theta<1\), has limiting cost \(\mathcal C(\alpha^\theta)=\theta\mathcal C(\alpha)<s\), so its recovered laws are eventually feasible at budget \(s\varepsilon^2\). Taking \(\varepsilon\downarrow0\) and then \(\theta\uparrow1\) in \eqref{eq:legacy-B.44} gives the lower bound \(F(s)\). Conversely, maximizing laws along a subsequence realizing the risk limsup admit, by the preceding compactness and risk-limit arguments, a limiting kernel of cost at most \(s\) with the same scaled risk, giving the upper bound. At zero the reference expansion applies directly. Finally, \eqref{eq:sharp-budget-continuity} gives
\begin{equation}
0\le D_\varepsilon(\rho_2)-D_\varepsilon(\rho_1)\le K(\rho_2-\rho_1).
\tag{A.48}\label{eq:legacy-B.50}
\end{equation}
Here \(K=(2\kappa_B+\kappa_H)/p\) bounds \(\kappa_\varepsilon\) for \(0<\varepsilon<1\). The rescaled profiles therefore share a Lipschitz constant, so pointwise convergence extends by a finite fine grid to the compact-uniform convergence \eqref{eq:legacy-4.4}.
\end{proof}

\subsection{A.6 Strict increase and critical-budget inversion}\label{app:strict-crossing-proof}\nopagebreak[4]
\label{rem:kernel-mixtures}
Finite cost does not ensure stationary mass: the kernel
\(\alpha_{\xi,u}=\delta_{\eta e^{-u^2}e_1}\), \(\eta>0\), moves every source and has
\(I=\eta\gamma_1g\sqrt\pi<\infty\).
Moreover, \(\mathcal V\) need not be concave along mixtures. To see this, choose a compact patch \(E\Subset\Omega\cap\{H\ge h>0\}\) and write
\(R=E\times[-h/4,-h/8]\), \(m_R=\nu(R)>0\), and
\(B_R=\int_R(u+H)\,d\nu>0\).
For \(|\tau|\le h/16\), the sources in \(R\) satisfy \(u\le-h/8<\tau\) and \(u+H\ge3h/4>\tau\). Since \(H\le\Delta_I\), the shift \(-2\Delta_Ie_1\) places both rescaled losses below \(-h/16\). Thus the mixture
\(\alpha^\theta=(1-\theta1_R)\delta_0+\theta1_R\delta_{-2\Delta_Ie_1}\)
has
\(T_0(\tau)=0\) and \(T_1(\tau)=\theta(m_R\tau-B_R)\).
Since \(I(\alpha^\theta)=2\Delta_I\gamma_1\theta m_R\to0\), the bound \(g\tau^2/2\le2\kappa_BI\) following \eqref{eq:legacy-4.2} places every minimizer in this window for sufficiently small \(\theta>0\). Minimizing the two quadratic objectives gives
\[
\mathcal V(\alpha^\theta)=-J/p+\theta B_R/p+\theta^2m_R^2/(2gp).
\]
Here minimizing \(G_1\) lowers its value by a quadratic term, which enters \(\mathcal V\) with the opposite sign because the second minimum is subtracted. The resulting positive quadratic term rules out a concavity argument for mixtures, so the replacement proof compares the two minimized threshold objectives directly.


\begin{proof}[Proof of Theorem~\ref{thm:strict-crossing}]\label{v21-crossing-proof}\label{appendix-c-strict-improvement-for-arbitrary-kernels-and-inversion-of-the-critical-budgets}
We apply the reflection from \mbox{\hyperref[sec:reflection-construction]{Appendix~A.5}} to an existing kernel, so its contribution on the replacement bands must also be controlled. Fix \(M\), retain the same patches \(E_\pm\), and choose \([-T,T]\) to contain every minimizing threshold for \(I\le M+1\). With \(a_T=\|H\|_\infty+T\), place unit-width bands beyond this window at
\(E=(E_+\times[U,U+1])\cup(E_-\times[-U-1,-U])\), \(U>a_T\).

By \eqref{eq:legacy-B.46}--\eqref{eq:legacy-B.47} with \(R=1\), the reflection \(d^N=-2ue_1\) has moment \(I_N=4\gamma_1m(U+1/2)\) and hinge contribution \(N_c(\tau)=-c\Delta_E\) throughout \([-T,T]\). For any old kernel with \(I\le M\), denote its moment and hinge contribution on \(E\) by \(I_O\le M\) and \(O_c(\tau)\). Subtracting the common baseline \(k_O=O_0(0)\) leaves \(R_c(\tau)=O_c(\tau)-k_O\), whose integrand is bounded by \(a_T\) and can be nonzero only if \(|d|_\gamma\ge(U-a_T)/\kappa_B\). Thus the finite displacement moment gives the uniform estimate
\begin{equation}
\sup_{|\tau|\le T}|R_c(\tau)|\le\frac{a_T\kappa_BM}{U-a_T}. \tag{A.49}\label{eq:legacy-C.2}
\end{equation}
The reflected gain \(\Delta_E\) does not depend on the band location, so we can fix \(U\), depending only on \(M\) and the model, with \(I_N\ge1\) and \(\sup|R_c|\le\Delta_E/4\).

For \(0<\delta\le1\), set \(\theta=\delta/(2I_N)\le1/2\), and replace the conditional kernel on \(E\) by \((1-\theta)\alpha+\theta\delta_{d^N}\), leaving it unchanged elsewhere. The resulting probability kernel \(\widehat\alpha_\delta=\alpha^\theta\) satisfies \(I(\widehat\alpha)=I(\alpha)+\theta(I_N-I_O)\) and \(j_i^{\widehat\alpha}=j_i^\alpha+\theta(j_i^N-j_i^O)\). Projection contraction bounds the old and new projected norms by \(I_O\) and \(I_N\), respectively, hence
\begin{equation}
\mathcal C(\widehat\alpha)-\mathcal C(\alpha)
\le\theta(I_N-I_O)+\theta(I_N+I_O)=\delta. \tag{A.50}\label{eq:legacy-C.3}
\end{equation}
All original and modified minimizers remain in \([-T,T]\), since \(I(\widehat\alpha)\le M+\delta/2\). On this common interval, \(G_c^{\widehat\alpha}(\tau)=G_c^\alpha(\tau)+\theta[-c\Delta_E-k_O-R_c(\tau)]\). The constant \(k_O\) cancels from the difference of minima, leaving the reflection gain with an error at most \(2\theta\sup|R_c|\). Consequently \(\mathcal V(\widehat\alpha)-\mathcal V(\alpha)\ge\theta\Delta_E/(2p)=c_M\delta\), where \(c_M=\Delta_E/(4pI_N)>0\).

Keeping these bands fixed as \(\delta\downarrow0\), we have \(\int|d|_\gamma|\alpha^\theta-\alpha|(dd)d\nu\le\theta(I_O+I_N)\to0\), implying convergence of the first moment, projected fluxes, and total cost. Since the objectives depend affinely on \(\theta\) on a common compact threshold interval, Danskin's formula \cite[Proposition~A.3.2(a)]{bertsekas2009}, with \(\mathcal M_c=\arg\min G_c^\alpha\), gives
\begin{equation}
\left.\frac d{d\theta}\mathcal V(\alpha^\theta)\right|_{0+}
=\frac{\Delta_E+\max_{\mathcal M_1}R_1-\max_{\mathcal M_0}R_0}{p}\ge\frac{\Delta_E}{2p}>0. \tag{A.51}\label{eq:legacy-C.4}
\end{equation}
Rescaling by \(\theta=\delta/(2I_N)\) gives the stated right derivative in \(\delta\).

Applying \eqref{eq:legacy-5.1} with \(M=S_0\) to a maximizer at budget \(s\), and subdividing if \(t-s>1\), proves \eqref{eq:legacy-5.2}. Together with continuity, \(F(0)<0\), and \eqref{eq:legacy-B.48}, this gives the unique zero \(s^*\). Its maximizing kernel uses cost \(s^*\), since smaller cost would give nonnegative risk below the zero. Local improvement produces positive values at costs decreasing to \(s^*\), while none can occur at \(s^*\) itself, proving the three characterizations and attainment assertions in \eqref{eq:legacy-5.3}. Finally, for \(0\le s_-<s^*<s_+\), compact-uniform convergence gives \(D_\varepsilon(s_-\varepsilon^2)<0<D_\varepsilon(s_+\varepsilon^2)\) for small \(\varepsilon\). Monotonicity in the budget places both critical budgets in \([s_-\varepsilon^2,s_+\varepsilon^2]\); letting \(s_\pm\to s^*\) proves \eqref{eq:legacy-5.4}.
\end{proof}

\subsection{A.7 Boundary sensitivity}\label{app:cost-proofs}

\leavevmode\phantomsection\label{l2.-the-baseline-minimizing-set-and-the-right-difference-quotient}
\begin{proof}[Proof of the boundary expansion in Theorem~\ref{prop:boundary-sensitivity}] We first justify minimization over the boundary equality kernels. On the measures \(\Pi_\alpha\), the coordinate functional \(A_i(\alpha)=\int|d_i|d\Pi_\alpha+\|j_i\|_1\) is lower semicontinuous for vague convergence by \mbox{\hyperref[v21-limit-proof]{Appendix~A.5}}. On \(\{\mathcal V=0\}\), the efficiency inequality gives \(\mathcal C_0\ge\gamma J\), so \(\mathcal C_0=\gamma J\) is equivalent to \(\mathcal C_0\le\gamma J\). Hence
\[
\mathcal Z_0=\{\alpha:\mathcal V(\alpha)=0,\ \mathcal C_0(\alpha)\le\gamma J\}.
\]
By compactness and continuity of \(\mathcal V\) on cost sublevels (Appendix~A.5), \(\mathcal Z_0\) is compact. It is nonempty because it contains the background comparison kernel \(\alpha^{(3)}\) from \eqref{eq:legacy-D.1}. The lower semicontinuous functional \(A_3\) therefore attains its minimum \(a_{3,\star}=A_3(\alpha_\star)\). Since \(j_3=0\) on \(\mathcal Z_0\), this minimum is the least background displacement moment, and comparison with \(\alpha^{(3)}\) gives \(0\le\beta a_{3,\star}\le J\). For \(C_\eta=\mathcal C_{\boldsymbol\gamma(\eta)}\),
\begin{equation}
C_\eta=C_0+\eta\beta\gamma A_3.
\tag{A.52}\label{eq:legacy-H.8}
\end{equation}
If \(\alpha_\eta\) minimizes the cost at zero risk difference, comparison with \(\alpha_\star\) gives
\begin{equation}
\gamma J\le C_0(\alpha_\eta)\le C_\eta(\alpha_\eta)
=s^*(\boldsymbol\gamma(\eta))
\le\gamma J+\eta\beta\gamma a_{3,\star}.
\tag{A.53}\label{eq:legacy-H.9}
\end{equation}
For \(0<\eta\le1\), \eqref{eq:legacy-H.9} places every minimizer in the compact sublevel \(\{C_0\le2\gamma J\}\). Along any vaguely convergent subsequence as \(\eta\downarrow0\), continuity of \(\mathcal V\) gives \(\mathcal V(\bar\alpha)=0\), and lower semicontinuity gives \(C_0(\bar\alpha)\le\gamma J\). These two properties place \(\bar\alpha\) in \(\mathcal Z_0\). Using the added background cost in \eqref{eq:legacy-H.8}, we obtain
\begin{equation}
\beta\gamma A_3(\alpha_\eta)
\le\frac{s^*(\boldsymbol\gamma(\eta))-\gamma J}{\eta}
\le\beta\gamma a_{3,\star}.
\tag{A.54}\label{eq:legacy-H.10}
\end{equation}
Along any vaguely convergent subsequence of minimizers, lower semicontinuity shows that the lower limit of the left side of \eqref{eq:legacy-H.10} is at least \(\beta\gamma A_3(\bar\alpha)\ge\beta\gamma a_{3,\star}\). The quotient therefore converges to \(\beta\gamma a_{3,\star}\), and every vague cluster point minimizes \(A_3\) on \(\mathcal Z_0\), without requiring uniqueness. Division by \(\gamma J\) gives the expansion with \(a_+\in[0,1]\). \end{proof}

\subsection{A.8 Equality kernels and continuous flow duality}\label{app:flow-proof}

\leavevmode\phantomsection\label{i.2.-from-arbitrary-randomized-kernels-in-mathcal-z_0-to-linear-fields-and-flows}
\textit{Proof of the flow representation in Theorem~\ref{thm:projected-flow}.} We reduce the background-cost minimization over \(\mathcal Z_0\) to a flow on \(D=\Omega\cap\{x<y\}\) by integrating out the normal coordinate and then identifying the two exchanged surface halves. The sign condition in Lemma~\ref{lem:efficiency-bound} is essential: on each half, all nonzero coordinate displacements have the sign of \(\zeta\), so integration in \(u\) introduces no cancellation in the absolute background displacement. For \(\alpha\in\mathcal Z_0\), set \(m_i(x,y)=\int_{\mathbb R}\int d_i\,\alpha_{\xi,u}(dd)du\) and \(\mathbf m=(m_1,m_2,\beta m_3)\). On \(D\), \(\mathbf m_i\ge0\) and \(\sum_i \mathbf m_i=w\); on its exchanged half, the signs reverse and the sum is \(-w\). Thus \(\beta\int|d_3|d\Pi_\alpha=\int_\Omega|\mathbf m_3|d\sigma\). To express the allocation using shares on \(D\) alone, we average the field with its exchanged counterpart:
\begin{equation}
(\mathscr S\mathbf m)(x,y)=(-\mathbf m_2(y,x),-\mathbf m_1(y,x),-\mathbf m_3(y,x)),
\qquad \bar{\mathbf m}=(\mathbf m+\mathscr S\mathbf m)/2.
\tag{A.55}\label{eq:legacy-I.9}
\end{equation}
The exchange interchanges the first two zero projections and preserves \(\xi_3=(v-x-y)/\beta\), so symmetrization preserves all three constraints. It also preserves \(\int|\mathbf m_3|d\sigma\), because the two background components have the same sign on each half. The shares \(a=\bar{\mathbf m}_1/w\), \(b=\bar{\mathbf m}_2/w\), \(c=\bar{\mathbf m}_3/w\) on \(D\) are therefore measurable, nonnegative and sum to one. On the exchanged half the components are \(-wb(y,x),-wa(y,x),-wc(y,x)\). To recover the flow constraint, we test either of the first two projections against a bounded measurable \(\varphi\) and obtain \(\int_D\varphi(x)a\,d\mu=\int_D\varphi(y)b\,d\mu\). This feasible flow has background amount
\begin{equation}
\beta\int|d_3|d\Pi_\alpha
=2\int_D(1-a-b)d\mu
=J-2\int_D(a+b)d\mu.
\tag{A.56}\label{eq:legacy-I.11}
\end{equation}

\paragraph{Lifting flows to probability kernels}\label{lem:v21-flow-lift}
Conversely, for feasible shares set \(\boldsymbol\theta=(a,b,1-a-b)\) on \(D\) and use \((b(y,x),a(y,x),1-a(y,x)-b(y,x))\) on the exchanged half. Since \(w>0\) on \(D\), the shares also define a simplex-valued allocation \(\sigma\)-almost everywhere. The kernel
\begin{equation}
\alpha_{\xi,u}=\delta_{\zeta(\theta_1,\theta_2,\theta_3/\beta)}
\tag{A.57}\label{eq:legacy-I.12}
\end{equation}
has finite first moment: its activity is confined to \(|u|\le\|H\|_\infty\), where \(|d|_1\le\max\{1,\beta^{-1}\}\|H\|_\infty\). Feasibility gives the first two zero projections, and exchange symmetry gives the third. Since its normal displacement is \(\zeta\) at minimum unit cost, Lemma~\ref{lem:efficiency-bound} places it in \(\mathcal Z_0\). Formula~\eqref{eq:legacy-I.11} then identifies the least background amount as \(J-2\mathfrak F_\mu\).

\paragraph{Continuous duality and the cut representation}
We enforce equality of the projected measures by testing their difference against a continuous function on \(I\). The allocation set \(\mathcal K=\{(a,b)\in L^\infty(\mu)^2:a,b\ge0,\ a+b\le1\}\) is convex and weak-star compact, being a weak-star closed subset of the product of two unit balls. For \(\varphi\in C(I)\), the corresponding Lagrangian is
\[
\mathscr H(a,b;\varphi)
=\int_D[(1-\varphi(x))a+(1+\varphi(y))b]\,d\mu.
\]
Since \(\varphi(x),\varphi(y)\in L^1(\mu)\), the Lagrangian \(\mathscr H\) is affine and weak-star continuous in \((a,b)\), and affine and continuous in \(\varphi\) for the uniform norm. Because continuous tests determine finite signed measures on \(I\), the conservation constraint defines a weak-star closed feasible subset of \(\mathcal K\) on which the flow maximum is attained. If the projected difference is nonzero, scaling a separating continuous test gives \(\inf_\varphi\mathscr H=-\infty\); for a feasible allocation, the potential terms cancel and the infimum equals \(\int_D(a+b)d\mu\). Sion's minimax theorem \cite[Theorem 4.2]{sion1958} applied to the real-valued \(\mathscr H\) on \(\mathcal K\times C(I)\) therefore yields
\[
\mathfrak F_\mu
=\sup_{(a,b)\in\mathcal K}\inf_{\varphi\in C(I)}\mathscr H(a,b;\varphi)
=\inf_{\varphi\in C(I)}\max_{(a,b)\in\mathcal K}\mathscr H(a,b;\varphi).
\]
For \(\varphi\in C(I)\), the Borel sets
\[
\begin{aligned}
E_a&=\{1-\varphi(x)>\max(0,1+\varphi(y))\},\\
E_b&=\{1+\varphi(y)\ge\max(0,1-\varphi(x))\}
      \cap\{1+\varphi(y)>0\}
\end{aligned}
\]
are disjoint. Their indicators \((a,b)=(1_{E_a},1_{E_b})\) belong to \(\mathcal K\) and attain the pointwise maximum, giving
\begin{equation}
D_\mu(\varphi)=\int_D\max\{0,1-\varphi(x),1+\varphi(y)\}\,d\mu.
\tag{A.58}\label{eq:legacy-I.13}
\end{equation}
For every bounded measurable potential, full projected equality still gives the weak-duality bound \(\int_D(a+b)d\mu\le D_\mu(\varphi)\). Enlarging the class of potentials therefore preserves the infimum obtained above. Truncation to \([-1,1]\) cannot increase the integrand, which proves \eqref{eq:legacy-6.8}.

To obtain a cut formula, we must reduce general bounded potentials to three-valued ones without changing the dual infimum. We do this by expressing the objective as an average over signed level sets. Let \(h_t(z)=\operatorname{sgn}(z)1_{\{|z|\ge t\}}\), \(0<t<1\). Because each \(h_t\) is nondecreasing, the signed layer-cake identity \(z=\int_0^1h_t(z)\,dt\) on \([-1,1]\) commutes with the maximum of two arguments. Since \(\max\{0,1+z_1,1+z_2\}=1+\max(z_1,z_2)\) on that square, applying the identity with \(z_1=-\varphi(x)\), \(z_2=\varphi(y)\) and \(\psi_t=1_{\{\varphi\ge t\}}-1_{\{\varphi\le-t\}}\) yields
\begin{equation}
D_\mu(\varphi)=\int_0^1D_\mu(\psi_t)\,dt.
\tag{A.59}\label{eq:legacy-I.23}
\end{equation}
Thus \(\inf_{0<t<1}D_\mu(\psi_t)\le D_\mu(\varphi)\), so the dual infimum is unchanged when restricted to \(\psi=1_U-1_V\) with \(U\cap V=\varnothing\). For these three-valued potentials, the integrand in \eqref{eq:legacy-I.13} has the explicit form
\[
\max\{0,1-\psi(x),1+\psi(y)\}
=1-1_{U\times V}(x,y)+1_{(V\times I)\cup(I\times U)}(x,y).
\]
Integrating over \(D\) and using \(\mu(D)=J/2\) gives \(D_\mu(\psi)=J/2-\mu(D\cap(U\times V))+\mu(D\cap((V\times I)\cup(I\times U)))\). Substitution into \(a_+=1-2\mathfrak F_\mu/J\) proves \eqref{eq:legacy-6.9}.

It remains to verify the strict flow bounds. For positivity, a two-step transfer \(K_1\to K_2\to K_3\) allows us to balance both projected measures on the middle interval \(K_2\). Choose three disjoint intervals \(K_1<K_2<K_3\) sufficiently close to the interior diagonal point from \hyperref[geometry:interior-diagonal]{Section~2.3} that the closed rectangles \(R_{12}=K_1\times K_2\), \(R_{23}=K_2\times K_3\) lie in \(D\). With \(\ell_{K,i}=|K_i|\), positivity and continuity of \(w_\sigma\) allow \(\kappa>0\) small enough that both densities
\[
bw_\sigma=\frac{\kappa}{\ell_{K,1}\ell_{K,2}}1_{R_{12}},
\qquad aw_\sigma=\frac{\kappa}{\ell_{K,2}\ell_{K,3}}1_{R_{23}},
\]
are bounded by \(w_\sigma\). Their disjoint supports give \(a+b\le1\), and both projected densities equal \(\kappa1_{K_2}/\ell_{K,2}\), producing a feasible flow of value \(2\kappa>0\). At the other endpoint, \(\mathfrak F_\mu=J/2\) would lift through \eqref{eq:legacy-I.12} to an equality kernel with no background displacement. Raising \(c_3\) above \(\gamma\) would leave its cost at \(\gamma J\), contrary to Theorem~\ref{thm:zero-premium}. Hence \(0<\mathfrak F_\mu<J/2\), and \eqref{eq:legacy-I.11} gives \(a_+\in(0,1)\). \unskip\nobreak\hfill\proofbox


\keepproofstart{8\baselineskip}
\section{Appendix B. Proofs for the analytical illustrations}\label{app:technical-quantitative}
\subsection{B.1 The projection-potential estimate}\label{app:potential-proof}
\begin{proof}[Proof of Proposition~\ref{prop:projection-potential-bound}]
Fix any minimizing threshold \(\tau\in\arg\min G_1\). With \(w=b\cdot d\), \(\Delta_+(u,w)=(u+w)_+-u_+\), and \(B_h(u,w)=\Delta_+(u,w)-\Delta_+(u+h,w)\), evaluating \(G_0\) at zero and \(G_1\) at \(\tau\) gives
\begin{equation}
J+p\mathcal V+\frac g2\tau^2
\le\int B_{H-\tau}(u,w)\,d\Pi_\alpha.
\tag{B.1}\label{eq:legacy-J.3}
\end{equation}
Boundedness of the potentials and finiteness of the first moment permit Fubini's theorem, giving
\[
\mathcal C\ge\int\left(|d|_\gamma+\sum_i\phi_i(\xi_i)d_i\right)d\Pi_\alpha,
\qquad
\inf_{b\cdot d=w}\left(|d|_\gamma+\sum_i\phi_i d_i\right)
=k_+w_++k_-(-w)_+.
\]
For \(w\ge0\), the positive parts of \(b_id_i\) sum to at least \(w\), and concentrating displacement in a coordinate of minimum effective cost attains \(k_+w\). Reversing signs gives \(k_-(-w)\) for \(w\le0\). It remains to maximize
\(P_h(u)=\sup_w\{\lambda B_h(u,w)-k_+w_+-k_-(-w)_+\}\).
If \(h<0\), then \(B_h(u,w)=r_h(u+w)-r_h(u)\) for the increasing ramp
\(r_h(t)=\min\{\max\{t,0\},|h|\}\), so only positive \(w\) can improve on zero. With \(k=k_+\), we have \(P_h=0\) when \(\lambda\le k\). When \(\lambda>k\), the nonzero pieces are
\(P_h(u)=[(\lambda-k)|h|+ku]_+\) for \(u<0\) and
\(P_h(u)=(\lambda-k)(|h|-u)\) for \(0\le u<|h|\).
Integrating these two triangles and reflecting when \(h>0\) gives
\[
\int_{\mathbb R}\sup_w\{\lambda B_h(u,w)-k_+w_+-k_-(-w)_+\}\,du
=\frac{\lambda h^2}{2}\left(\frac{\lambda}{k(\xi,h)}-1\right)_+.
\]
The value is zero for \(h=0\). Since \(H\) is bounded and \(k_\pm\) are bounded away from zero, integrating this estimate over \(\Omega\) in \eqref{eq:legacy-J.3} proves \eqref{eq:legacy-J.6}.
\end{proof}

\subsection{B.2 The uniform midpoint-potential bound}\label{app:uniform-bound-proof}

To establish the two threshold branches used for the lower bound in Corollary~\ref{prop:uniform-variable-cost}, we normalize costs and potentials by \(\gamma_0\), substitute the midpoint potentials of \hyperref[strict-premium-despite-equal-efficiency]{Section~5.1}, and take \(\lambda=3/2\) in these normalized units. We first evaluate the penalty on the same-side regions \(\{x<y<1\}\) and \(\{1<x<y\}\), intersected with \(\Omega\), together with their exchanged images. On these regions, both effective directional costs equal \(1/2\). Their total section mass is \(1/8\), their \(H^2\)-moment is \(1/48\), and their \(H\)-moment is zero. Their contribution to the integral of \((H-\tau)^2\) is consequently \(1/48+\tau^2/8\).

On the cross-region, \(k(\xi,H)=3/2=\lambda\): the direction that increases the paired hinge gain has sign opposite to \(H\), so the penalty vanishes at \(\tau=0\). For \(0<|\tau|\le1\), a penalty appears only where the threshold shift reverses the sign of \(H-\tau\). In that triangular strip, the coordinates \(h=|H|\) and \(t=x+y\) satisfy \(0<h<|\tau|\), \(2-h<t<2+h\), and \(d\sigma=dt\,dh/16\). Thus the additional integral is
\[
\int_0^{|\tau|}(|\tau|-h)^2\frac h8\,dh=\frac{\tau^4}{96}.
\]
The penalty multiplier in \eqref{eq:legacy-J.6} is \(\frac\lambda2\{\lambda/(1/2)-1\}=3/2\). Substitution, with \(g=3/8\) and \(J=5/32\), yields
\[
\frac{\mathcal C}{\gamma_0}
\ge\frac32\left(\frac5{32}+\frac12\mathcal V+\frac3{16}\tau^2\right)
-\frac32\left(\frac1{48}+\frac{\tau^2}{8}+\frac{\tau^4}{96}\right)
=\frac34\mathcal V+\frac{13}{64}+\frac{\tau^2(6-\tau^2)}{64}.
\]
For \(|\tau|\ge1\), the nonnegative terms in \eqref{eq:legacy-E.5} instead give \(\mathcal C/\gamma_0\ge\mathcal V/2+J+g/2=\mathcal V/2+11/32\). These estimates cover every minimizing threshold.

\subsection{B.3 Proof of Proposition~\ref{prop:dependence-family}}\label{app:dependence-family-proof}

\begin{proof}
To use all the same-side mass, we draw matching projections from the disjoint regions
\(E_A=\{x<1<y,\ x+y>2\}\cap D\) and
\(E_B=\{x<1<y,\ x+y<2\}\cap D\). Define the required and available section masses by
\[
\begin{aligned}
d_{\vartheta,-}(q)&=\int 1_{D_-}(x,q)w_{\sigma,\vartheta}(x,q)\,dx,&
Q_{\vartheta,A}(q)&=\int1_{E_A}(q,y)w_{\sigma,\vartheta}(q,y)\,dy
&& (q<1),\\
d_{\vartheta,+}(q)&=\int1_{D_+}(q,y)w_{\sigma,\vartheta}(q,y)\,dy,&
Q_{\vartheta,B}(q)&=\int1_{E_B}(x,q)w_{\sigma,\vartheta}(x,q)\,dx
&& (q>1).
\end{aligned}
\]
Integration on the lower half, with \(w=1-q\in[0,1/2]\), gives
\[
Q_{\vartheta,A}(q)-d_{\vartheta,-}(q)
=w\left\{\frac3{16}-\frac{3w}{16}+\frac{w^2}{48}
+\vartheta\left(-\frac1{16}+\frac{w}{8}
-\frac{w^3}{16}-\frac{w^4}{30}\right)\right\}.
\]
For every \(|\vartheta|<1/2\), the polynomial in braces is strictly decreasing on \([0,1/2]\) with a positive endpoint value. Since \(d_{\vartheta,-}=0\) below \(q=1/2\), we obtain \(Q_{\vartheta,A}\ge d_{\vartheta,-}\) throughout the lower half. Reflecting about \(q=1\) gives \(Q_{\vartheta,B}\ge d_{\vartheta,+}\) on the upper half. The disjoint matching allocation is therefore
\begin{equation}
(a,b)=
\begin{cases}
(0,1),&D_-,\\
(1,0),&D_+,\\
(d_{\vartheta,-}(x)/Q_{\vartheta,A}(x),0),&E_A,\\
(0,d_{\vartheta,+}(y)/Q_{\vartheta,B}(y)),&E_B.
\end{cases}\tag{B.2}\label{eq:legacy-K.7}
\end{equation}
The capacity inequalities and disjointness of \(E_A,E_B\) give \(a,b\ge0\), \(a+b\le1\), and
\begin{equation}
\int a(q,y)w_{\sigma,\vartheta}(q,y)\,dy
=\int b(x,q)w_{\sigma,\vartheta}(x,q)\,dx
=
\begin{cases}
d_{\vartheta,-}(q),&q<1,\\
d_{\vartheta,+}(q),&q>1,
\end{cases}
\quad\text{for almost every }q.
\tag{B.3}\label{eq:legacy-K.8}
\end{equation}
This projected equality gives a feasible flow of value \(2(M_{\vartheta,-}+M_{\vartheta,+})\), attained dually by \(\varphi=1_{[0,1)}-1_{(1,2]}\), whose integrand is \(2\) on \(D_-\cup D_+\) and zero elsewhere. Theorem~\ref{thm:projected-flow} then gives the background share. We evaluate \(J_\vartheta\) using exchange symmetry and reflection about \(t=2\), in coordinates \(t=x+y\), \(h=y-x>0\) with Jacobian \(1/2\). For the same-side masses, we use \(x\)-sections of \(D_-\) and \(u=2y-1\in(0,1)\). These integrations give
\[
\begin{aligned}
J_\vartheta
&=\frac18\int_1^2\int_0^t
h^2\{1-\vartheta(t-2)^2\}\,dh\,dt
=\frac{75-14\vartheta}{480},\\
M_{\vartheta,-}=M_{\vartheta,+}
&=\frac1{960}\int_0^1
u^3\{10+\vartheta(-10+5u-u^2)\}\,du
=\frac{3-2\vartheta}{1152}.
\end{aligned}
\]
\end{proof}


\keepproofstart{10\baselineskip}
\section*{Declarations}
\addcontentsline{toc}{section}{Declarations}
\subsection*{Use of artificial intelligence}
In accordance with the SIAM Editorial Policy on Artificial Intelligence, the author reports using DeepSeek to assist with brainstorming, literature review, language editing, and adversarial stress-testing. The author independently verified all mathematical arguments, derivations, and cited references, and assumes full responsibility for all content.

\subsection*{Funding and competing interests}
The author received no external funding for this research and declares no competing interests.

\subsection*{Data and code availability}\label{computational-supplement}
No external datasets were used in this study. The reproducible Python scripts used to verify the analytical polynomial integrals and to generate Figures 1--4 are available in the accompanying code archive.

\begin{thebibliography}{99}
\bibitem{artzner1999}
\textsc{P. Artzner, F. Delbaen, J.-M. Eber, and D. Heath}, \emph{Coherent measures of risk}, Math. Finance, 9(3) (1999), pp. 203--228, \href{https://doi.org/10.1111/1467-9965.00068}{doi:10.1111/1467-9965.00068}.

\bibitem{bartl2021}
\textsc{D. Bartl, S. Drapeau, J. Obłój, and J. Wiesel}, \emph{Sensitivity analysis of Wasserstein distributionally robust optimization problems}, Proc. R. Soc. A, 477(2256) (2021), article 20210176, \href{https://doi.org/10.1098/rspa.2021.0176}{doi:10.1098/rspa.2021.0176}.

\bibitem{bartleckstein2026}
\textsc{D. Bartl and S. Eckstein}, \emph{Optimal nonparametric estimation of the expected shortfall risk}, preprint (2026), to appear in SIAM J. Financial Math., \href{https://arxiv.org/abs/2405.00357v2}{arXiv:2405.00357v2}.

\bibitem{bartlwiesel2023}
\textsc{D. Bartl and J. Wiesel}, \emph{Sensitivity of multiperiod optimization problems with respect to the adapted Wasserstein distance}, SIAM J. Financial Math., 14(2) (2023), pp. 704--720, \href{https://doi.org/10.1137/22M1537746}{doi:10.1137/22M1537746}.

\bibitem{bernard2014}
\textsc{C. Bernard, X. Jiang, and R. Wang}, \emph{Risk aggregation with dependence uncertainty}, Insur. Math. Econ., 54 (2014), pp. 93--108, \href{https://doi.org/10.1016/j.insmatheco.2013.11.005}{doi:10.1016/j.insmatheco.2013.11.005}.

\bibitem{bertsekas2009}
\textsc{D. P. Bertsekas}, \emph{Convex Optimization Theory}, Athena Scientific, Belmont, MA, 2009, \href{https://web.mit.edu/dimitrib/www/Convex_Theory_Entire_Book.pdf}{author's online edition}.

\bibitem{blanchet2019}
\textsc{J. Blanchet and K. Murthy}, \emph{Quantifying distributional model risk via optimal transport}, Math. Oper. Res., 44(2) (2019), pp. 565--600, \href{https://doi.org/10.1287/moor.2018.0936}{doi:10.1287/moor.2018.0936}.

\bibitem{brandtner2018}
\textsc{M. Brandtner}, \emph{Expected shortfall, spectral risk measures, and the aggravating effect of background risk, or: risk vulnerability and the problem of subadditivity}, J. Bank. Financ., 89 (2018), pp. 138--149, \href{https://doi.org/10.1016/j.jbankfin.2018.02.002}{doi:10.1016/j.jbankfin.2018.02.002}.

\bibitem{chenfurmanlin2026}
\textsc{J. Chen, E. Furman, and X. S. Lin}, \emph{Static marginal expected shortfall: Systemic risk measurement under dependence uncertainty}, Insur. Math. Econ., 130 (2026), article 103285, \href{https://doi.org/10.1016/j.insmatheco.2026.103285}{doi:10.1016/j.insmatheco.2026.103285}.

\bibitem{chen2022}
\textsc{Y. Chen, P. Liu, Y. Liu, and R. Wang}, \emph{Ordering and inequalities for mixtures on risk aggregation}, Math. Finance, 32(1) (2022), pp. 421--451, \href{https://doi.org/10.1111/mafi.12323}{doi:10.1111/mafi.12323}.

\bibitem{dalmaso1993}
\textsc{G. Dal Maso}, \emph{An Introduction to $\Gamma$-Convergence}, Progress in Nonlinear Differential Equations and Their Applications 8, Birkh\"auser, Boston, 1993, \href{https://doi.org/10.1007/978-1-4612-0327-8}{doi:10.1007/978-1-4612-0327-8}.

\bibitem{devecchinendelstreicher2026}
\textsc{C. De Vecchi, M. Nendel, and J. Streicher}, \emph{Upper comonotonicity and risk aggregation under dependence uncertainty}, Math. Finance, 36(1) (2026), pp. 118--139, \href{https://doi.org/10.1111/mafi.70008}{doi:10.1111/mafi.70008}.

\bibitem{eckstein2020}
\textsc{S. Eckstein, M. Kupper, and M. Pohl}, \emph{Robust risk aggregation with neural networks}, Math. Finance, 30(4) (2020), pp. 1229--1272, \href{https://doi.org/10.1111/mafi.12280}{doi:10.1111/mafi.12280}.

\bibitem{eckstein2021multiasset}
\textsc{S. Eckstein, G. Guo, T. Lim, and J. Obłój}, \emph{Robust pricing and hedging of options on multiple assets and its numerics}, SIAM J. Financial Math., 12(1) (2021), pp. 158--188, \href{https://doi.org/10.1137/19M1286256}{doi:10.1137/19M1286256}.

\bibitem{embrechts2013}
\textsc{P. Embrechts, G. Puccetti, and L. R\"uschendorf}, \emph{Model uncertainty and VaR aggregation}, J. Bank. Financ., 37(8) (2013), pp. 2750--2764, \href{https://doi.org/10.1016/j.jbankfin.2013.03.014}{doi:10.1016/j.jbankfin.2013.03.014}.

\bibitem{embrechts2015}
\textsc{P. Embrechts, B. Wang, and R. Wang}, \emph{Aggregation-robustness and model uncertainty of regulatory risk measures}, Finance Stoch., 19(4) (2015), pp. 763--790, \href{https://doi.org/10.1007/s00780-015-0273-z}{doi:10.1007/s00780-015-0273-z}.

\bibitem{fiechtner2026}
\textsc{L.-B. Fiechtner and J. Blanchet}, \emph{Wasserstein distributionally robust regret optimization}, preprint, 2025, revised 2026, \href{https://arxiv.org/abs/2504.10796v5}{arXiv:2504.10796v5}.

\bibitem{gao2017}
\textsc{R. Gao and A. J. Kleywegt}, \emph{Data-driven robust optimization with known marginal distributions}, working paper, 2017, \url{https://faculty.mccombs.utexas.edu/rui.gao/copula.pdf}.

\bibitem{gao2023}
\textsc{R. Gao and A. J. Kleywegt}, \emph{Distributionally robust stochastic optimization with Wasserstein distance}, Math. Oper. Res., 48(2) (2023), pp. 603--655, \href{https://doi.org/10.1287/moor.2022.1275}{doi:10.1287/moor.2022.1275}.

\bibitem{korman2015}
\textsc{J. Korman, R. J. McCann, and C. Seis}, \emph{Dual potentials for capacity constrained optimal transport}, Calc. Var. Partial Differential Equations, 54 (2015), pp. 573--584, \href{https://doi.org/10.1007/s00526-014-0795-9}{doi:10.1007/s00526-014-0795-9}.

\bibitem{lam2018}
\textsc{H. Lam}, \emph{Sensitivity to serial dependency of input processes: A robust approach}, Manage. Sci., 64(3) (2018), pp. 1311--1327, \href{https://doi.org/10.1287/mnsc.2016.2667}{doi:10.1287/mnsc.2016.2667}.

\bibitem{li2018partial}
\textsc{L. Li, H. Shao, R. Wang, and J. Yang}, \emph{Worst-case Range Value-at-Risk with partial information}, SIAM J. Financial Math., 9(1) (2018), pp. 190--218, \href{https://doi.org/10.1137/17M1126138}{doi:10.1137/17M1126138}.

\bibitem{liu2026}
\textsc{P. Liu and Y. Liu}, \emph{Portfolio diversification and concentration under dependence uncertainty: A majorization approach}, preprint, 2026, \href{https://arxiv.org/abs/2609.04496v1}{arXiv:2609.04496v1}.

\bibitem{martin2011}
\textsc{R. J. Martin}, \emph{Saddlepoint methods in portfolio theory}, in The Oxford Handbook of Credit Derivatives, A. Lipton and A. Rennie, eds., Oxford University Press, Oxford (2011), pp. 533--569, \href{https://doi.org/10.1093/oxfordhb/9780199546787.013.0015}{doi:10.1093/oxfordhb/9780199546787.013.0015}.

\bibitem{mcneil2015}
\textsc{A. J. McNeil, R. Frey, and P. Embrechts}, \emph{Quantitative Risk Management: Concepts, Techniques and Tools}, revised ed., Princeton University Press, Princeton, NJ, 2015.

\bibitem{namkoong2026}
\textsc{H. Namkoong}, \emph{Local sensitivity under transport restrictions}, preprint, 2026, \href{https://arxiv.org/abs/2606.04276v1}{arXiv:2606.04276v1}.

\bibitem{otakimura2021}
\textsc{S. Ota and M. Kimura}, \emph{Effective estimation algorithm for parameters of multivariate Farlie--Gumbel--Morgenstern copula}, Jpn. J. Stat. Data Sci., 4 (2021), pp. 1049--1078, \href{https://doi.org/10.1007/s42081-021-00118-y}{doi:10.1007/s42081-021-00118-y}.

\bibitem{pesenti2019}
\textsc{S. M. Pesenti, P. Millossovich, and A. Tsanakas}, \emph{Reverse sensitivity testing: What does it take to break the model?}, European J. Oper. Res., 274(2) (2019), pp. 654--670, \href{https://doi.org/10.1016/j.ejor.2018.10.003}{doi:10.1016/j.ejor.2018.10.003}.

\bibitem{pesenti2022}
\textsc{S. M. Pesenti}, \emph{Reverse sensitivity analysis for risk modelling}, Risks, 10(7) (2022), article 141, \href{https://doi.org/10.3390/risks10070141}{doi:10.3390/risks10070141}.

\bibitem{pesentijaimungal2023}
\textsc{S. M. Pesenti and S. Jaimungal}, \emph{Portfolio optimization within a Wasserstein ball}, SIAM J. Financial Math., 14(4) (2023), pp. 1175--1214, \href{https://doi.org/10.1137/22M1496803}{doi:10.1137/22M1496803}.

\bibitem{pflugpohl2018}
\textsc{G. Ch. Pflug and M. Pohl}, \emph{A review on ambiguity in stochastic portfolio optimization}, Set-Valued Var. Anal., 26(4) (2018), pp. 733--757, \href{https://doi.org/10.1007/s11228-017-0458-z}{doi:10.1007/s11228-017-0458-z}.

\bibitem{rockafellar2002}
\textsc{R. T. Rockafellar and S. Uryasev}, \emph{Conditional value-at-risk for general loss distributions}, J. Bank. Financ., 26(7) (2002), pp. 1443--1471, \href{https://doi.org/10.1016/S0378-4266(02)00271-6}{doi:10.1016/S0378-4266(02)00271-6}.

\bibitem{ruschendorf2013}
\textsc{L. R\"uschendorf}, \emph{Mathematical Risk Analysis: Dependence, Risk Bounds, Optimal Allocations and Portfolios}, Springer Series in Operations Research and Financial Engineering, Springer, Berlin, Heidelberg, 2013, \href{https://doi.org/10.1007/978-3-642-33590-7}{doi:10.1007/978-3-642-33590-7}.

\bibitem{santambrogio2015}
\textsc{F. Santambrogio}, \emph{Optimal Transport for Applied Mathematicians: Calculus of Variations, PDEs, and Modeling}, Progress in Nonlinear Differential Equations and Their Applications 87, Birkh\"auser, Cham, 2015, \href{https://doi.org/10.1007/978-3-319-20828-2}{doi:10.1007/978-3-319-20828-2}.

\bibitem{sauldubois2026}
\textsc{N. Sauldubois}, \emph{Higher-order sensitivity analysis of Wasserstein and adapted Wasserstein distributionally robust optimization problems}, preprint, 2026, \href{https://arxiv.org/abs/2609.25345v1}{arXiv:2609.25345v1}.

\bibitem{scaillet2004}
\textsc{O. Scaillet}, \emph{Nonparametric estimation and sensitivity analysis of expected shortfall}, Math. Finance, 14(1) (2004), pp. 115--129, \href{https://doi.org/10.1111/j.0960-1627.2004.00184.x}{doi:10.1111/j.0960-1627.2004.00184.x}.

\bibitem{sion1958}
\textsc{M. Sion}, \emph{On general minimax theorems}, Pacific J. Math., 8(1) (1958), pp. 171--176, \href{https://doi.org/10.2140/pjm.1958.8.171}{doi:10.2140/\allowbreak pjm.1958.8.171}.

\bibitem{tasche2001}
\textsc{D. Tasche}, \emph{Conditional expectation as quantile derivative}, preprint, 2001, \href{https://arxiv.org/abs/math/0104190v1}{arXiv:math/0104190v1}.

\bibitem{vallender1974}
\textsc{S. S. Vallender}, \emph{Calculation of the Wasserstein distance between probability distributions on the line}, Theory Probab. Appl., 18(4) (1974), pp. 784--786, \href{https://doi.org/10.1137/1118101}{doi:10.1137/1118101}.

\bibitem{villani2003}
\textsc{C. Villani}, \emph{Topics in Optimal Transportation}, Graduate Studies in Mathematics 58, American Mathematical Society, Providence, RI, 2003.

\bibitem{wang2015superadditive}
\textsc{R. Wang, V. Bignozzi, and A. Tsanakas}, \emph{How superadditive can a risk measure be?}, SIAM J. Financial Math., 6(1) (2015), pp. 776--803, \href{https://doi.org/10.1137/140981046}{doi:10.1137/140981046}.

\end{thebibliography}

\end{document}
